% Alternate 1, derived directly from the original exam-solution pair.
% Build from this directory; all dependencies remain in this academic-year tree.
\makeatletter
\def\input@path{{../../../preamble/}}
\makeatother
\input{preamble_exams.tex}
\input{highlights/highlights.tex}
% \GTenHighlightsfalse % Uncomment to disable instructional highlights.
\input{highlights/specific_highlights.tex}
\CoordinatedTermTwoPalette
\tikzset{gten marker/.append style={inner xsep=.6mm,inner ysep=.3mm}}
\usepackage{varwidth}
\usepackage{needspace}
% Natural cell widths, bounded locally so long descriptions can wrap.
\newcommand{\CreditDescription}[1]{\begin{varwidth}[t]{.72\linewidth}\raggedright#1\end{varwidth}}
% Match the reference practice graph treatments; the toggle also controls ink.
\colorlet{keyMassColor}{blue!75!black}
\colorlet{keyCDFColor}{cyan!65!blue!80!black}
\colorlet{keyMassFill}{GTenHLProbability!32!white}
\colorlet{keyGuideColor}{keyMassColor!35}
\ifGTenHighlights\else
  \colorlet{keyMassColor}{black}
  \colorlet{keyCDFColor}{black}
  \colorlet{keyMassFill}{black!15}
  \colorlet{keyGuideColor}{gray}
\fi
\linespread{0.98}
\setlength{\parskip}{2pt}
\renewcommand{\arraystretch}{1.0}
\setlength{\emergencystretch}{2em}
\pgfplotsset{keyaxis/.style={axis lines=left,scaled ticks=false,
 xlabel={$\GTenHLValue{x}$ (\GTenHLValue{USD})},
 xticklabel={ $\GTenHLValue{\pgfmathprintnumber{\tick}}$},grid=major,grid style={gray!15},
 tick label style={font=\small},label style={font=\small},clip=false}}
\pgfplotsset{keyCombined/.style={width=.84\linewidth,height=4cm,scale only axis,
every axis plot/.append style={forget plot},axis lines=left,x axis line style={->},ymin=-.035,ymax=1.13,
ytick={0,.25,.5,.75,1},grid=major,
grid style={gray!15},tick label style={font=\small},
label style={font=\small},clip=false,
legend style={at={(.5,-.34)},anchor=north,legend columns=2,
font=\footnotesize,draw=none,cells={anchor=west}},
},keyMass/.style={keyMassColor,thick},
keyPhantom/.style={keyGuideColor,densely dashed,line width=.9pt},
keyStep/.style={keyCDFColor,thick},
keyIncluded/.style={only marks,mark=*,keyCDFColor,mark size=1.1pt},
keyExcluded/.style={only marks,mark=*,keyCDFColor,
mark options={fill=white},mark size=1.1pt}}
% Arguments: supported value, exact mass, cumulative height immediately before it.
% Only a vertical translation is applied: upper height = lower height + mass.
% Coordinates are saved inside their own axes, then joined after both axes end.
\usetikzlibrary{backgrounds,calc}
\tikzset{keyCorrespondence/.style={gray!70,line width=.25pt,-{Stealth[length=3.5pt,width=3pt]}}}
\newcommand{\keyCorrespondence}[1]{%
\begingroup
% Maximum sideways bow in physical units, independent of the axis scales.
\def\keyCorrespondenceBow{3mm}%
% Quadratic Bezier with midpoint control shifted by twice the bow;
% its equivalent cubic controls lie at 1/3 and 2/3, shifted by 4/3 of it.
\pgfmathsetlengthmacro{\keyControlOffset}{4/3*\keyCorrespondenceBow}%
% Paint behind the stems and markers, even though both axes are already defined.
\begin{pgfonlayer}{background}
\draw[keyCorrespondence] (keyMassMid-#1)
  .. controls ([xshift=\keyControlOffset]$(keyMassMid-#1)!1/3!(keyShadowMid-#1)$)
          and ([xshift=\keyControlOffset]$(keyMassMid-#1)!2/3!(keyShadowMid-#1)$)
  .. (keyShadowMid-#1);
\end{pgfonlayer}%
\endgroup}
\newcommand{\keyPhantom}[3]{%
\coordinate (keyShadowMid-#1) at (axis cs:#1,{#3+(#2)/2});
\addplot[keyPhantom] coordinates {(#1,#3)(#1,{#3+#2})};
\addplot[only marks,mark=*,keyGuideColor,mark size=1.8pt]
coordinates {(#1,{#3+#2})};}
\newcommand{\keyMass}[2]{%
\coordinate (keyMassMid-#1) at (axis cs:#1,{(#2)/2});
\addplot[keyMass] coordinates {(#1,0)(#1,#2)};
\addplot[only marks,mark=*,keyMassColor,mark size=1.5pt]
coordinates {(#1,#2)};}

\newcommand{\keyConnectionReading}{%
\textbf{Reading the comparison.} Both plotting areas use identical horizontal
limits, support ticks, numerical spacing, widths and vertical scales.
Each pale dashed shadow translates its lower \GTenHLProbability{PMF} stem upward
by $\HLCumulativeProbability{F_X(}\GTenHLValue{a^-}\HLCumulativeProbability{)}$ without changing its length:
$\HLCumulativeProbability{F_X(}\GTenHLValue{a}\HLCumulativeProbability{)}-\HLCumulativeProbability{F_X(}\GTenHLValue{a^-}\HLCumulativeProbability{)}=\GTenHLProbability{p_X(}\GTenHLValue{a}\GTenHLProbability{)}$.
Thin curved arrows join the stem and shadow midpoints, including the first
support value. They indicate correspondence. Stems and shadows show isolated
masses; the filled \HLCumulativeProbability{CDF} endpoint gives its value and
the open endpoint excludes the lower height.}

\title{Global Economics Exam Solutions\\Grade 10 --- Term 2 --- C1 and C2}
\author{}
\date{}
\makeatletter
\renewcommand{\@maketitle}{\begin{center}{\Large\bfseries\@title\par}\end{center}}
\renewcommand\subsection{\@startsection{subsection}{2}{\z@}{-1.4ex plus -.3ex}{.7ex}{\normalfont\large\bfseries}}
\makeatother
\hypersetup{pdftitle={Grade 10 Term 2 C1 and C2 Exam Solutions},pdfauthor={Global Economics}}
\begin{document}
\maketitle
\small
\setlength{\topsep}{3pt}
\setcounter{tocdepth}{1}
\phantomsection
\label{toc}
\tableofcontents
\section*{Learning objective}
Model discrete monetary outcomes in economic and financial contexts; derive and verify probability mass functions (PMFs), and construct and interpret cumulative distribution functions (CDFs) with tables and graphs.
\section*{Criteria assessment}
This key reproduces each alternate question before its worked solution. Questions 1--2 assess C1; Questions 3--6 assess C2. Show reasoning and exact probabilities. All contextual rules are stated in ordinary language.

This exam assesses C1 and C2. A \textbf{minimum of 4 out of 6 points for C1 and 4 out of 14 points for C2} is required to pass. The overall available total is 20 points. Show all relevant calculations and justify your answers. Partial credit may be awarded for correct procedures despite arithmetic errors; answers without demonstrated procedures receive no credit. A 1-point penalty per criterion applies if the student's name is missing or the assessment is submitted late.

\textbf{Independent passing thresholds: C1 requires at least 4 out of 6 points; C2 requires at least 4 out of 14 points. Points cannot transfer between criteria.} The overall available total is 20 points. An overall total alone does not establish a pass in both criteria. These are the passing thresholds for this assessment; no cumulative policy is specified.
\section*{Evaluated criteria}
\begin{description}
\item[C1] Characterises discrete random variables in economic and financial contexts by identifying their support, probability mass function and possible outcomes.
\item[C2] Interprets PMFs and CDFs by verifying their defining properties and using them to calculate point, interval, cumulative, and complementary probabilities in financial contexts.
\end{description}

\begin{center}
\begin{tabular}{|c|c|c|}
	\hline\textbf{\textnumero} & \textbf{C1} & \textbf{C2} \\ \hline
	\ref*{sec:q1} & 3 & --- \\\hline
	\ref*{sec:q2} & 3 & --- \\\hline
	\ref*{sec:q3} & --- & 2 \\\hline
	\ref*{sec:q4} & --- & 2 \\\hline
	\ref*{sec:q5} & --- & 4 \\\hline
	\ref*{sec:q6} & --- & 6 \\\hline
	\textbf{Total} & \textbf{6} & \textbf{14} \\\hline
\end{tabular}
\end{center}

\textbf{Marking policy.} Question and criterion totals are fixed by the brief. Task totals match the student exam; detailed partial-credit breakdowns below are proposed. Award stated fractions for demonstrated components, including valid alternative methods. Allow consistent follow-through after one arithmetic error when subsequent reasoning is correct. In C1 Task 2 (Questions 1--2) and C2 Task 1 (Questions 3--6), positivity (nonnegativity) and normalization checks carry no separate points; omitting these checks does not reduce the task score. Solving the normalization equation to derive probabilities in Question 5 remains assessed, as do its relationship checks. CDF verification retains its stated credit and must be demonstrated. Cap each task, question and criterion at its allocation.


\clearpage
\section*{Highlight legend}
\phantomsection\label{tab:highlight-legend}
\begin{center}
\begin{tabular}{@{}lll@{}}\toprule
\textbf{Symbol or notation} & \textbf{Concept} & \textbf{Brief meaning}\\\midrule
$\GTenHLOutcome{\omega}$ & \GTenHLOutcome{Outcome} & One complete result or category.\\
$\GTenHLValue{X},\GTenHLValue{x},\GTenHLValue{S_X}$ & \GTenHLValue{Variable and support} & Values and cutoffs in \GTenHLValue{USD}.\\
$\GTenHLSecond{N},\GTenHLSecond{Y}$ & \GTenHLSecond{Other quantities} & Counts; adjusted payout in USD.\\
$\GTenHLProbability{P},\GTenHLProbability{p_X(}\GTenHLValue{x}\GTenHLProbability{)}$ & \GTenHLProbability{Probability / PMF} & Unitless likelihood or point mass.\\
$\HLCumulativeProbability{F_X(}\GTenHLValue{x}\HLCumulativeProbability{)}$ & \HLCumulativeProbability{Cumulative probability / CDF} & Mass at or below a cutoff.\\
$\GTenHLParameter{a}$ & \GTenHLParameter{Unknown parameter} & Baseline mass to determine in Question 5.\\
\bottomrule\end{tabular}
\end{center}

\sectionScored{Meal reservation contribution}{C1-\sml{3}}{sec:q1}

A meal vendor handles three meal reservations. The three reservations are independent. Each reservation is either kept or cancelled, two mutually exclusive and exhaustive outcomes. The keep probabilities for reservations 1, 2 and 3 are 0.2, 0.5 and 0.8, respectively. Each kept reservation contributes 7 USD and each cancellation contributes no money. A fixed handling cost of 5 USD is then deducted once from the combined contribution. There are no other receipts or costs. Record the net contribution for the three meal reservations.

\Needspace{8\baselineskip}
\subsection*{Task 1 (a) --- question and worked solution (proposed allocation) \score{1}}
\label{sol:q1t1}

\noindent\textbf{\boldmath Define the random variable and its units, and identify its support.}\par\nopagebreak

Let $\GTenHLValue{X}$ be the \GTenHLValue{net contribution} from the three meal reservations, in \GTenHLValue{USD}. If $\GTenHLSecond{N}$ is the number of kept reservations, the monetary rule is
\[
\GTenHLValue{X}=\GTenHLValue{7}\GTenHLSecond{N}-\GTenHLValue{5}.
\]
Since $\GTenHLSecond{N}\in\{0,1,2,3\}$, the distinct possible values are $7(0)-5=-5$, $7(1)-5=2$, $7(2)-5=9$ and $7(3)-5=16$ USD. Therefore
\[
\GTenHLValue{S_X}=\{\GTenHLValue{-5},\GTenHLValue{2},\GTenHLValue{9},\GTenHLValue{16}\}\ \GTenHLValue{USD}.
\]
A negative net contribution means the handling cost exceeds receipts.

\par\smallskip
\noindent\begin{minipage}{\linewidth}\centering
\textbf{Proposed partial credit.}\par\smallskip
\begin{tabular}{@{}l r@{}}\toprule
\CreditDescription{\textbf{Assessed step or criterion}} & \textbf{Points}\\\midrule
\CreditDescription{Contextual definition of $X$, monetary rule and units} & 0.5\\
\CreditDescription{Correct, complete support} & 0.5\\
\midrule
\CreditDescription{\textbf{Total}} & \textbf{1}\\
\bottomrule\end{tabular}
\end{minipage}\par\smallskip

\par\vspace{3mm}\secInicio[][sec:q1]

\clearpage
\subsection*{Task 2 (b) --- question and worked solution (proposed allocation) \score{1}}
\label{sol:q1t2}

\noindent\textbf{\boldmath List every complete outcome, preserving order where relevant; map each outcome to its probability and monetary value; and construct the grouped PMF.}\par\nopagebreak

Write $\GTenHLOutcome{K}$ for \GTenHLOutcome{kept} and $\GTenHLOutcome{C}$ for \GTenHLOutcome{cancelled}. Each triple is in reservation 1, 2, 3 order. Before calculating products, use the complementary outcomes:
\[
\begin{aligned}
P(C_1)&=1-P(K_1)=1-0.2=\GTenHLProbability{0.8},\\
P(C_2)&=1-P(K_2)=1-0.5=\GTenHLProbability{0.5},\\
P(C_3)&=1-P(K_3)=1-0.8=\GTenHLProbability{0.2}.
\end{aligned}
\]
There are $2^3=8$ ordered complete outcomes. Independence permits multiplication using each reservation's own probabilities; the trials have different keep probabilities.
\begin{center}
\begin{tabular}{@{}cccc@{}}\toprule
Ordered outcome & Probability calculation & Exact $P(\{\omega\})$ & $\GTenHLValue{X}(\GTenHLOutcome{\omega})$ (USD)\\
\midrule
$\GTenHLOutcome{\mathrm{KKK}}$ & $\GTenHLProbability{0.2}\cdot\GTenHLProbability{0.5}\cdot\GTenHLProbability{0.8}$ & $\GTenHLProbability{2/25}=\GTenHLProbability{0.08}=\GTenHLProbability{8\%}$ & $\GTenHLValue{16}$\\
$\GTenHLOutcome{\mathrm{KKC}}$ & $\GTenHLProbability{0.2}\cdot\GTenHLProbability{0.5}\cdot\GTenHLProbability{0.2}$ & $\GTenHLProbability{1/50}=\GTenHLProbability{0.02}=\GTenHLProbability{2\%}$ & $\GTenHLValue{9}$\\
$\GTenHLOutcome{\mathrm{KCK}}$ & $\GTenHLProbability{0.2}\cdot\GTenHLProbability{0.5}\cdot\GTenHLProbability{0.8}$ & $\GTenHLProbability{2/25}=\GTenHLProbability{0.08}=\GTenHLProbability{8\%}$ & $\GTenHLValue{9}$\\
$\GTenHLOutcome{\mathrm{KCC}}$ & $\GTenHLProbability{0.2}\cdot\GTenHLProbability{0.5}\cdot\GTenHLProbability{0.2}$ & $\GTenHLProbability{1/50}=\GTenHLProbability{0.02}=\GTenHLProbability{2\%}$ & $\GTenHLValue{2}$\\
$\GTenHLOutcome{\mathrm{CKK}}$ & $\GTenHLProbability{0.8}\cdot\GTenHLProbability{0.5}\cdot\GTenHLProbability{0.8}$ & $\GTenHLProbability{8/25}=\GTenHLProbability{0.32}=\GTenHLProbability{32\%}$ & $\GTenHLValue{9}$\\
$\GTenHLOutcome{\mathrm{CKC}}$ & $\GTenHLProbability{0.8}\cdot\GTenHLProbability{0.5}\cdot\GTenHLProbability{0.2}$ & $\GTenHLProbability{2/25}=\GTenHLProbability{0.08}=\GTenHLProbability{8\%}$ & $\GTenHLValue{2}$\\
$\GTenHLOutcome{\mathrm{CCK}}$ & $\GTenHLProbability{0.8}\cdot\GTenHLProbability{0.5}\cdot\GTenHLProbability{0.8}$ & $\GTenHLProbability{8/25}=\GTenHLProbability{0.32}=\GTenHLProbability{32\%}$ & $\GTenHLValue{2}$\\
$\GTenHLOutcome{\mathrm{CCC}}$ & $\GTenHLProbability{0.8}\cdot\GTenHLProbability{0.5}\cdot\GTenHLProbability{0.2}$ & $\GTenHLProbability{2/25}=\GTenHLProbability{0.08}=\GTenHLProbability{8\%}$ & $\GTenHLValue{-5}$\\
\bottomrule\end{tabular}
\end{center}
Exhaustiveness check: $(0.2+0.8)(0.5+0.5)(0.8+0.2)=1$; the eight probabilities sum to 1.


\par\vspace{3mm}\secInicio[][sec:q1]

\clearpage
\textbf{Question 1 --- Task 2 (b) continued: grouped \GTenHLProbability{PMF}}

For one complete outcome $\GTenHLOutcome{\omega}$, $\GTenHLProbability{P(\{}\GTenHLOutcome{\omega}\GTenHLProbability{\})}$ is its probability. The grouped mass adds the probabilities of all outcomes with the same monetary value:
\[
\GTenHLProbability{p_X(}\GTenHLValue{x}\GTenHLProbability{)}
=\sum_{\omega:\,X(\omega)=x}\GTenHLProbability{P(\{}\GTenHLOutcome{\omega}\GTenHLProbability{\})}.
\]
\begin{center}
\begin{tabular}{@{}c l l@{}}\toprule
$\GTenHLValue{x}$ (\GTenHLValue{USD}) & Complete outcomes & Grouped $\GTenHLProbability{p_X(}\GTenHLValue{x}\GTenHLProbability{)}$\\\midrule
$\GTenHLValue{-5}$ & $\GTenHLOutcome{\mathrm{CCC}}$ & $\begin{aligned}[t]\GTenHLProbability{8/100}&=\GTenHLProbability{8/100}\\&=\GTenHLProbability{0.08}=\GTenHLProbability{8\%}\end{aligned}$\\[5pt]
$\GTenHLValue{2}$ & $\GTenHLOutcome{\mathrm{KCC}},\GTenHLOutcome{\mathrm{CKC}},\GTenHLOutcome{\mathrm{CCK}}$ & $\begin{aligned}[t]\GTenHLProbability{2/100}+\GTenHLProbability{8/100}+\GTenHLProbability{32/100}&=\GTenHLProbability{42/100}\\&=\GTenHLProbability{0.42}=\GTenHLProbability{42\%}\end{aligned}$\\[5pt]
$\GTenHLValue{9}$ & $\GTenHLOutcome{\mathrm{KKC}},\GTenHLOutcome{\mathrm{KCK}},\GTenHLOutcome{\mathrm{CKK}}$ & $\begin{aligned}[t]\GTenHLProbability{2/100}+\GTenHLProbability{8/100}+\GTenHLProbability{32/100}&=\GTenHLProbability{42/100}\\&=\GTenHLProbability{0.42}=\GTenHLProbability{42\%}\end{aligned}$\\[5pt]
$\GTenHLValue{16}$ & $\GTenHLOutcome{\mathrm{KKK}}$ & $\begin{aligned}[t]\GTenHLProbability{8/100}&=\GTenHLProbability{8/100}\\&=\GTenHLProbability{0.08}=\GTenHLProbability{8\%}\end{aligned}$\\[5pt]
\bottomrule\end{tabular}
\end{center}
The grouped PMF is the last column.
\Needspace{5\baselineskip}\textbf{PMF checks.}
\begin{itemize}
\item Nonnegativity: all listed masses are positive.
\item Normalization: $(\GTenHLProbability{8}+\GTenHLProbability{42}+\GTenHLProbability{42}+\GTenHLProbability{8})/\GTenHLProbability{100}=\GTenHLProbability{1}$.
\item Outside the \GTenHLValue{support}: the \GTenHLProbability{PMF} is zero elsewhere.
\end{itemize}

\textbf{Scoring note.} These explanatory positivity and normalization checks carry no separate points; omitting them does not reduce the task score.

\par\smallskip
\noindent\begin{minipage}{\linewidth}\centering
\textbf{Proposed partial credit.}\par\smallskip
\begin{tabular}{@{}l r@{}}\toprule
\CreditDescription{\textbf{Assessed step or criterion}} & \textbf{Points}\\\midrule
\CreditDescription{Complete ordered outcomes} & 0.25\\
\CreditDescription{Correct outcome probability calculations} & 0.25\\
\CreditDescription{Correct monetary mapping for every outcome} & 0.25\\
\CreditDescription{Correct grouping and grouped PMF masses} & 0.25\\
\midrule
\CreditDescription{\textbf{Total}} & \textbf{1}\\
\bottomrule\end{tabular}
\end{minipage}\par\smallskip

\par\vspace{3mm}\secInicio[][sec:q1]

\Needspace{8\baselineskip}
\subsection*{Task 3 (c) --- question and worked solution (proposed allocation) \score{1}}
\label{sol:q1t3}

\noindent\textbf{\boldmath Calculate and interpret the probability that the net contribution is exactly 9 USD.}\par\nopagebreak

From the grouped \GTenHLProbability{PMF}, $P(\GTenHLValue{X}=\GTenHLValue{9})=\GTenHLProbability{p_X(}\GTenHLValue{9}\GTenHLProbability{)}=\GTenHLProbability{\frac{21}{50}}=\GTenHLProbability{0.42}=\GTenHLProbability{42\%}$. The vendor earns a \GTenHLValue{net contribution} of exactly \GTenHLValue{9 USD} with probability \GTenHLProbability{42\%}; this happens exactly when two reservations are kept and one is cancelled.

\par\smallskip
\noindent\begin{minipage}{\linewidth}\centering
\textbf{Proposed partial credit.}\par\smallskip
\begin{tabular}{@{}l r@{}}\toprule
\CreditDescription{\textbf{Assessed step or criterion}} & \textbf{Points}\\\midrule
\CreditDescription{Selecting the correct mass at $X=9$} & 0.5\\
\CreditDescription{The contextual interpretation} & 0.5\\
\midrule
\CreditDescription{\textbf{Total}} & \textbf{1}\\
\bottomrule\end{tabular}
\end{minipage}\par\smallskip

Question total: 3 points for C1.

\par\vspace{3mm}\secInicio[][sec:q1]

\clearpage
\sectionScored{Two distributor promotional credits}{C1-\sml{3}}{sec:q2}

A bookshop receives one promotional credit from each of two distributors. The distributors award credits independently. Each has exactly three mutually exclusive possible outcomes: no credit, a 4 USD credit, or a 9 USD credit. For the first distributor these probabilities are 0.3, 0.4 and 0.3, respectively. For the second distributor they are 0.5, 0.2 and 0.3, respectively. Add both credits; there are no other adjustments. Record the total credit, keeping first-distributor then second-distributor order.

\Needspace{8\baselineskip}
\subsection*{Task 1 (a) --- question and worked solution (proposed allocation) \score{1}}
\label{sol:q2t1}

\noindent\textbf{\boldmath Define the random variable and its units, and identify its support.}\par\nopagebreak

Let $\GTenHLValue{X}$ be the sum of the first and second distributors' promotional credits, in \GTenHLValue{USD}. The monetary rule adds the two credits. Each credit belongs to $\{0,4,9\}$ USD, so their possible sums are
\[
\begin{array}{lll}
0+0=0,&0+4=4,&0+9=9,\\
4+0=4,&4+4=8,&4+9=13,\\
9+0=9,&9+4=13,&9+9=18.
\end{array}
\]
Repeated sums appear only once in the support. Hence
\[
\GTenHLValue{S_X}=\{\GTenHLValue{0},\GTenHLValue{4},\GTenHLValue{8},\GTenHLValue{9},\GTenHLValue{13},\GTenHLValue{18}\}\ \GTenHLValue{USD}.
\]

\par\smallskip
\noindent\begin{minipage}{\linewidth}\centering
\textbf{Proposed partial credit.}\par\smallskip
\begin{tabular}{@{}l r@{}}\toprule
\CreditDescription{\textbf{Assessed step or criterion}} & \textbf{Points}\\\midrule
\CreditDescription{Contextual definition of $X$, monetary rule and units} & 0.5\\
\CreditDescription{Correct, complete support} & 0.5\\
\midrule
\CreditDescription{\textbf{Total}} & \textbf{1}\\
\bottomrule\end{tabular}
\end{minipage}\par\smallskip

\par\vspace{3mm}\secInicio[][sec:q2]

\clearpage
\subsection*{Task 2 (b) --- question and worked solution (proposed allocation) \score{1}}
\label{sol:q2t2}

\noindent\textbf{\boldmath List every complete outcome, preserving order where relevant; map each outcome to its probability and monetary value; and construct the grouped PMF.}\par\nopagebreak

Each distributor has three possible outcomes, giving $3^2=9$ ordered pairs. Write each pair in first-distributor then second-distributor order, in USD. Independence permits multiplication of the two marginal probabilities.
\begin{center}
\begin{tabular}{@{}cccc@{}}\toprule
Ordered pair (USD) & Probability calculation & Exact $P(\{\omega\})$ & $\GTenHLValue{X}(\GTenHLOutcome{\omega})$ (USD)\\\midrule
$(\GTenHLOutcome{0},\GTenHLOutcome{0})$ & $\GTenHLProbability{0.3}\cdot\GTenHLProbability{0.5}$ & $\GTenHLProbability{3/20}=\GTenHLProbability{0.15}=\GTenHLProbability{15\%}$ & $\GTenHLValue{0}$\\
$(\GTenHLOutcome{0},\GTenHLOutcome{4})$ & $\GTenHLProbability{0.3}\cdot\GTenHLProbability{0.2}$ & $\GTenHLProbability{3/50}=\GTenHLProbability{0.06}=\GTenHLProbability{6\%}$ & $\GTenHLValue{4}$\\
$(\GTenHLOutcome{0},\GTenHLOutcome{9})$ & $\GTenHLProbability{0.3}\cdot\GTenHLProbability{0.3}$ & $\GTenHLProbability{9/100}=\GTenHLProbability{0.09}=\GTenHLProbability{9\%}$ & $\GTenHLValue{9}$\\
$(\GTenHLOutcome{4},\GTenHLOutcome{0})$ & $\GTenHLProbability{0.4}\cdot\GTenHLProbability{0.5}$ & $\GTenHLProbability{1/5}=\GTenHLProbability{0.20}=\GTenHLProbability{20\%}$ & $\GTenHLValue{4}$\\
$(\GTenHLOutcome{4},\GTenHLOutcome{4})$ & $\GTenHLProbability{0.4}\cdot\GTenHLProbability{0.2}$ & $\GTenHLProbability{2/25}=\GTenHLProbability{0.08}=\GTenHLProbability{8\%}$ & $\GTenHLValue{8}$\\
$(\GTenHLOutcome{4},\GTenHLOutcome{9})$ & $\GTenHLProbability{0.4}\cdot\GTenHLProbability{0.3}$ & $\GTenHLProbability{3/25}=\GTenHLProbability{0.12}=\GTenHLProbability{12\%}$ & $\GTenHLValue{13}$\\
$(\GTenHLOutcome{9},\GTenHLOutcome{0})$ & $\GTenHLProbability{0.3}\cdot\GTenHLProbability{0.5}$ & $\GTenHLProbability{3/20}=\GTenHLProbability{0.15}=\GTenHLProbability{15\%}$ & $\GTenHLValue{9}$\\
$(\GTenHLOutcome{9},\GTenHLOutcome{4})$ & $\GTenHLProbability{0.3}\cdot\GTenHLProbability{0.2}$ & $\GTenHLProbability{3/50}=\GTenHLProbability{0.06}=\GTenHLProbability{6\%}$ & $\GTenHLValue{13}$\\
$(\GTenHLOutcome{9},\GTenHLOutcome{9})$ & $\GTenHLProbability{0.3}\cdot\GTenHLProbability{0.3}$ & $\GTenHLProbability{9/100}=\GTenHLProbability{0.09}=\GTenHLProbability{9\%}$ & $\GTenHLValue{18}$\\
\bottomrule\end{tabular}
\end{center}
Exhaustiveness check: $(0.3+0.4+0.3)(0.5+0.2+0.3)=1$. The ordered pairs are not equally likely.


\par\vspace{3mm}\secInicio[][sec:q2]

\clearpage
\textbf{Question 2 --- Task 2 (b) continued: grouped \GTenHLProbability{PMF}}

For one complete outcome $\GTenHLOutcome{\omega}$, $\GTenHLProbability{P(\{}\GTenHLOutcome{\omega}\GTenHLProbability{\})}$ is its probability. The grouped mass adds the probabilities of all outcomes with the same monetary value:
\[
\GTenHLProbability{p_X(}\GTenHLValue{x}\GTenHLProbability{)}
=\sum_{\omega:\,X(\omega)=x}\GTenHLProbability{P(\{}\GTenHLOutcome{\omega}\GTenHLProbability{\})}.
\]
\begin{center}
\begin{tabular}{@{}c l l@{}}\toprule
$\GTenHLValue{x}$ (\GTenHLValue{USD}) & Ordered outcomes & Grouped $\GTenHLProbability{p_X(}\GTenHLValue{x}\GTenHLProbability{)}$\\\midrule
$\GTenHLValue{0}$ & $(\GTenHLOutcome{0},\GTenHLOutcome{0})$ & $\begin{aligned}[t]\GTenHLProbability{3/20}&=\GTenHLProbability{3/20}\\&=\GTenHLProbability{0.15}=\GTenHLProbability{15\%}\end{aligned}$\\[5pt]
$\GTenHLValue{4}$ & $(\GTenHLOutcome{0},\GTenHLOutcome{4}),(\GTenHLOutcome{4},\GTenHLOutcome{0})$ & $\begin{aligned}[t]\GTenHLProbability{3/50}+\GTenHLProbability{1/5}&=\GTenHLProbability{13/50}\\&=\GTenHLProbability{0.26}=\GTenHLProbability{26\%}\end{aligned}$\\[5pt]
$\GTenHLValue{8}$ & $(\GTenHLOutcome{4},\GTenHLOutcome{4})$ & $\begin{aligned}[t]\GTenHLProbability{2/25}&=\GTenHLProbability{2/25}\\&=\GTenHLProbability{0.08}=\GTenHLProbability{8\%}\end{aligned}$\\[5pt]
$\GTenHLValue{9}$ & $(\GTenHLOutcome{0},\GTenHLOutcome{9}),(\GTenHLOutcome{9},\GTenHLOutcome{0})$ & $\begin{aligned}[t]\GTenHLProbability{9/100}+\GTenHLProbability{3/20}&=\GTenHLProbability{6/25}\\&=\GTenHLProbability{0.24}=\GTenHLProbability{24\%}\end{aligned}$\\[5pt]
$\GTenHLValue{13}$ & $(\GTenHLOutcome{4},\GTenHLOutcome{9}),(\GTenHLOutcome{9},\GTenHLOutcome{4})$ & $\begin{aligned}[t]\GTenHLProbability{3/25}+\GTenHLProbability{3/50}&=\GTenHLProbability{9/50}\\&=\GTenHLProbability{0.18}=\GTenHLProbability{18\%}\end{aligned}$\\[5pt]
$\GTenHLValue{18}$ & $(\GTenHLOutcome{9},\GTenHLOutcome{9})$ & $\begin{aligned}[t]\GTenHLProbability{9/100}&=\GTenHLProbability{9/100}\\&=\GTenHLProbability{0.09}=\GTenHLProbability{9\%}\end{aligned}$\\[5pt]
\bottomrule\end{tabular}
\end{center}
All masses are positive, their sum is $(15+26+8+24+18+9)/100=1$, and the PMF is zero elsewhere. Reversing a pair can preserve the monetary value while changing its probability.

\textbf{Scoring note.} These explanatory positivity and normalization checks carry no separate points; omitting them does not reduce the task score.

\par\smallskip
\noindent\begin{minipage}{\linewidth}\centering
\textbf{Proposed partial credit.}\par\smallskip
\begin{tabular}{@{}l r@{}}\toprule
\CreditDescription{\textbf{Assessed step or criterion}} & \textbf{Points}\\\midrule
\CreditDescription{Complete ordered outcomes} & 0.25\\
\CreditDescription{Correct outcome probability calculations} & 0.25\\
\CreditDescription{Correct monetary mapping for every outcome} & 0.25\\
\CreditDescription{Correct grouping and grouped PMF masses} & 0.25\\
\midrule
\CreditDescription{\textbf{Total}} & \textbf{1}\\
\bottomrule\end{tabular}
\end{minipage}\par\smallskip

\par\vspace{3mm}\secInicio[][sec:q2]

\Needspace{8\baselineskip}
\subsection*{Task 3 (c) --- question and worked solution (proposed allocation) \score{1}}
\label{sol:q2t3}

\noindent\textbf{\boldmath Calculate and interpret the probability that the total credit is exactly 13 USD.}\par\nopagebreak

$P(\GTenHLValue{X}=\GTenHLValue{13})=\GTenHLProbability{p_X(}\GTenHLValue{13}\GTenHLProbability{)}=\GTenHLProbability{3/25}+\GTenHLProbability{3/50}=\GTenHLProbability{\frac{9}{50}}=\GTenHLProbability{0.18}=\GTenHLProbability{18\%}$. The bookshop receives exactly 13 USD with probability \GTenHLProbability{18\%}; the qualifying pairs are $(4,9)$ and $(9,4)$.

\par\smallskip
\noindent\begin{minipage}{\linewidth}\centering
\textbf{Proposed partial credit.}\par\smallskip
\begin{tabular}{@{}l r@{}}\toprule
\CreditDescription{\textbf{Assessed step or criterion}} & \textbf{Points}\\\midrule
\CreditDescription{Selecting the correct mass at $X=13$} & 0.5\\
\CreditDescription{The contextual interpretation} & 0.5\\
\midrule
\CreditDescription{\textbf{Total}} & \textbf{1}\\
\bottomrule\end{tabular}
\end{minipage}\par\smallskip

Question total: 3 points for C1. Criterion C1 total: 6 points; independent pass: at least 4 points.

\par\vspace{3mm}\secInicio[][sec:q2]

\clearpage
\sectionScored{Produce crate delivery fees}{C2-\sml{2}}{sec:q3}

A farm cooperative sends one produce crate using exactly one of eight exhaustive, mutually exclusive delivery methods: walking courier, bicycle, cargo bicycle, motorcycle, car, van, refrigerated van or truck. Their probabilities are 0.15, 0.10, 0.10, 0.15, 0.20, 0.05, 0.15 and 0.10, respectively. The corresponding fixed handling fees are 0, 0, 6, 6, 10, 10, 21 and 21 USD, respectively. No other handling charges apply. Record the handling fee for the crate.

\Needspace{8\baselineskip}
\subsection*{Task 1 --- question and worked solution (proposed allocation) \score{0.5}}
\label{sol:q3t1}

\noindent\textbf{\boldmath Define the random variable and support; derive the grouped PMF; produce its table and a labeled graph with actual monetary spacing.}\par\nopagebreak

Let $\GTenHLValue{X}$ be the handling fee for one produce crate, in USD. Map all eight categories as follows.
\begin{center}
\begin{tabular}{ccc}\toprule
Category & Given probability & $\GTenHLValue{X}$ (USD)\\\midrule
\GTenHLOutcome{Walking courier} & $\GTenHLProbability{0.15}=\GTenHLProbability{3/20}$ & $\GTenHLValue{0}$\\
\GTenHLOutcome{Bicycle} & $\GTenHLProbability{0.10}=\GTenHLProbability{2/20}$ & $\GTenHLValue{0}$\\
\GTenHLOutcome{Cargo bicycle} & $\GTenHLProbability{0.10}=\GTenHLProbability{2/20}$ & $\GTenHLValue{6}$\\
\GTenHLOutcome{Motorcycle} & $\GTenHLProbability{0.15}=\GTenHLProbability{3/20}$ & $\GTenHLValue{6}$\\
\GTenHLOutcome{Car} & $\GTenHLProbability{0.20}=\GTenHLProbability{4/20}$ & $\GTenHLValue{10}$\\
\GTenHLOutcome{Van} & $\GTenHLProbability{0.05}=\GTenHLProbability{1/20}$ & $\GTenHLValue{10}$\\
\GTenHLOutcome{Refrigerated van} & $\GTenHLProbability{0.15}=\GTenHLProbability{3/20}$ & $\GTenHLValue{21}$\\
\GTenHLOutcome{Truck} & $\GTenHLProbability{0.10}=\GTenHLProbability{2/20}$ & $\GTenHLValue{21}$\\
\bottomrule\end{tabular}
\end{center}
There are eight distinct sample-space outcomes and four monetary values. Add probabilities of mutually exclusive outcomes with the same value:

$\GTenHLValue{S_X}=\{0,6,10,21\}$ USD.
\par\vspace{3mm}\secInicio[][sec:q3]
\clearpage\textbf{Question 3 --- Task 1 continued: grouped \GTenHLProbability{PMF} and graph}
\begin{center}
\begingroup\setlength{\tabcolsep}{5pt}\renewcommand{\arraystretch}{1.18}
\begin{tabular}{@{}p{.30\linewidth}c p{.55\linewidth}@{}}\toprule
Complete outcomes or categories & $\GTenHLValue{x}$ (\GTenHLValue{USD}) & Grouped probability: $\GTenHLProbability{p_X(}\GTenHLValue{x}\GTenHLProbability{)}$\\\midrule
\GTenHLOutcome{Walking courier}, \GTenHLOutcome{Bicycle} & $\GTenHLValue{0}$ & $\begin{aligned}[t]\GTenHLProbability{p_X(}\GTenHLValue{0}\GTenHLProbability{)}&=\GTenHLProbability{3/20}+\GTenHLProbability{2/20}=\GTenHLProbability{5/20}\\&=\GTenHLProbability{0.25}=\GTenHLProbability{25\%}\end{aligned}$\\[5pt]
\GTenHLOutcome{Cargo bicycle}, \GTenHLOutcome{Motorcycle} & $\GTenHLValue{6}$ & $\begin{aligned}[t]\GTenHLProbability{p_X(}\GTenHLValue{6}\GTenHLProbability{)}&=\GTenHLProbability{2/20}+\GTenHLProbability{3/20}=\GTenHLProbability{5/20}\\&=\GTenHLProbability{0.25}=\GTenHLProbability{25\%}\end{aligned}$\\[5pt]
\GTenHLOutcome{Car}, \GTenHLOutcome{Van} & $\GTenHLValue{10}$ & $\begin{aligned}[t]\GTenHLProbability{p_X(}\GTenHLValue{10}\GTenHLProbability{)}&=\GTenHLProbability{4/20}+\GTenHLProbability{1/20}=\GTenHLProbability{5/20}\\&=\GTenHLProbability{0.25}=\GTenHLProbability{25\%}\end{aligned}$\\[5pt]
\GTenHLOutcome{Refrigerated van}, \GTenHLOutcome{Truck} & $\GTenHLValue{21}$ & $\begin{aligned}[t]\GTenHLProbability{p_X(}\GTenHLValue{21}\GTenHLProbability{)}&=\GTenHLProbability{3/20}+\GTenHLProbability{2/20}=\GTenHLProbability{5/20}\\&=\GTenHLProbability{0.25}=\GTenHLProbability{25\%}\end{aligned}$\\[5pt]
\bottomrule\end{tabular}
\endgroup
\end{center}
\begin{center}\begin{tikzpicture}
\begin{axis}[keyaxis,width=.86\linewidth,height=5cm,xmin=-2,xmax=23,ymin=0,ymax=.65,xtick={0,6,10,21},ylabel={$\GTenHLProbability{p_X(}\GTenHLValue{x}\GTenHLProbability{)}$},yticklabel={$\GTenHLProbability{\pgfmathprintnumber{\tick}}$}]
\draw[fill=keyMassFill,draw=keyMassColor] (axis cs:-0.45,0) rectangle (axis cs:0.45,5/20);
\draw[fill=keyMassFill,draw=keyMassColor] (axis cs:5.55,0) rectangle (axis cs:6.45,5/20);
\draw[fill=keyMassFill,draw=keyMassColor] (axis cs:9.55,0) rectangle (axis cs:10.45,5/20);
\draw[fill=keyMassFill,draw=keyMassColor] (axis cs:20.55,0) rectangle (axis cs:21.45,5/20);
\end{axis}\end{tikzpicture}\end{center}\Needspace{5\baselineskip}\textbf{PMF checks.}\begin{itemize}\item Every mass is positive.\item Normalization: $(5+5+5+5)/20=1$.\item Set $\GTenHLProbability{p_X(}\GTenHLValue{x}\GTenHLProbability{)}=\GTenHLProbability{0}$ outside the support. Bars use actual monetary spacing; gaps carry no mass.\end{itemize}

\textbf{Scoring note.} These explanatory positivity and normalization checks carry no separate points; omitting them does not reduce the task score.

\par\smallskip
\noindent\begin{minipage}{\linewidth}\centering
\textbf{Proposed partial credit.}\par\smallskip
\begin{tabular}{@{}l r@{}}\toprule
\CreditDescription{\textbf{Assessed step or criterion}} & \textbf{Points}\\\midrule
\CreditDescription{Variable, units and support} & 0.125\\
\CreditDescription{Category mapping, grouping and the PMF table} & 0.25\\
\CreditDescription{The labeled, correctly spaced graph} & 0.125\\
\midrule
\CreditDescription{\textbf{Total}} & \textbf{0.5}\\
\bottomrule\end{tabular}
\end{minipage}\par\smallskip

\par\vspace{3mm}\secInicio[][sec:q3]

\clearpage
\textbf{Question 3 --- continued}\quad\hyperref[sec:q3]{Original description}

\Needspace{8\baselineskip}
\subsection*{Task 2 --- question and worked solution (proposed allocation) \score{1}}
\label{sol:q3t2}

\noindent\textbf{\boldmath Construct the complete CDF table, including both tails, and draw a labeled step graph. Use open and filled endpoints correctly and preserve actual monetary spacing.}\par\nopagebreak

For $\HLCumulativeProbability{F_X(}\GTenHLValue{x}\HLCumulativeProbability{)}=P(X\leq x)$, add masses in increasing monetary order.

The successive cumulative numerators are $5,10,15,20$, all over $20$.
\begin{center}
\begingroup
\setlength{\tabcolsep}{5pt}
\begin{tabular}{@{}r@{}c@{}c@{}c@{}llc@{}}\toprule
\multicolumn{5}{c}{Cutoff interval (USD)} & Calculations & $\HLCumulativeProbability{F_X(}\GTenHLValue{x}\HLCumulativeProbability{)}$\\\midrule
 & & $\GTenHLValue{x}$ & $<$ & $\GTenHLValue{0}$ & $\text{No included support values}$ & $\HLCumulativeProbability{0}$\\[5pt]
$\GTenHLValue{0}$ & $\leq$ & $\GTenHLValue{x}$ & $<$ & $\GTenHLValue{6}$ & $\begin{aligned}[t]
  &\GTenHLProbability{p_X(}\GTenHLValue{0}\GTenHLProbability{)}
  \\&= \GTenHLProbability{\frac{5}{20}}
\end{aligned}$ & $\HLCumulativeProbability{5/20}$\\[5pt]
$\GTenHLValue{6}$ & $\leq$ & $\GTenHLValue{x}$ & $<$ & $\GTenHLValue{10}$ & $\begin{aligned}[t]
  &\GTenHLProbability{p_X(}\GTenHLValue{0}\GTenHLProbability{)} + \GTenHLProbability{p_X(}\GTenHLValue{6}\GTenHLProbability{)}
  \\&= \GTenHLProbability{\frac{5}{20}} + \GTenHLProbability{\frac{5}{20}}
\end{aligned}$ & $\HLCumulativeProbability{10/20}$\\[5pt]
$\GTenHLValue{10}$ & $\leq$ & $\GTenHLValue{x}$ & $<$ & $\GTenHLValue{21}$ & $\begin{aligned}[t]
  &\GTenHLProbability{p_X(}\GTenHLValue{0}\GTenHLProbability{)} + \GTenHLProbability{p_X(}\GTenHLValue{6}\GTenHLProbability{)} + \GTenHLProbability{p_X(}\GTenHLValue{10}\GTenHLProbability{)}
  \\&= \GTenHLProbability{\frac{5}{20}} + \GTenHLProbability{\frac{5}{20}} + \GTenHLProbability{\frac{5}{20}}
\end{aligned}$ & $\HLCumulativeProbability{15/20}$\\[5pt]
$\GTenHLValue{21}$ & $\leq$ & $\GTenHLValue{x}$ &  &  & $\begin{aligned}[t]
  &\GTenHLProbability{p_X(}\GTenHLValue{0}\GTenHLProbability{)} + \GTenHLProbability{p_X(}\GTenHLValue{6}\GTenHLProbability{)} + \GTenHLProbability{p_X(}\GTenHLValue{10}\GTenHLProbability{)} + \GTenHLProbability{p_X(}\GTenHLValue{21}\GTenHLProbability{)}
  \\&= \GTenHLProbability{\frac{5}{20}} + \GTenHLProbability{\frac{5}{20}} + \GTenHLProbability{\frac{5}{20}} + \GTenHLProbability{\frac{5}{20}}
\end{aligned}$ & $\HLCumulativeProbability{1}$\\[5pt]
\bottomrule\end{tabular}
\endgroup
\end{center}
\begin{center}\begin{tikzpicture}
\begin{axis}[keyaxis,width=.86\linewidth,height=7cm,xmin=-2,xmax=23,ymin=-.07,ymax=1.1,xtick={0,6,10,21},ylabel={$\HLCumulativeProbability{F_X(}\GTenHLValue{x}\HLCumulativeProbability{)}$},yticklabel={$\HLCumulativeProbability{\pgfmathprintnumber{\tick}}$},ytick={0,.25,.5,.75,1}]
\addplot[keyCDFColor,thick,<-] coordinates {(-2,0)(0,0)};
\addplot[keyCDFColor,thick] coordinates {(0,5/20)(6,5/20)};
\addplot[keyCDFColor,thick] coordinates {(6,10/20)(10,10/20)};
\addplot[keyCDFColor,thick] coordinates {(10,15/20)(21,15/20)};
\addplot[keyCDFColor,thick,->] coordinates {(21,1)(23,1)};
\addplot[keyCDFColor,only marks,mark=*,mark size=1.6pt,mark options={fill=white}] coordinates {(0,0/20)(6,5/20)(10,10/20)(21,15/20)};
\addplot[keyCDFColor,only marks,mark=*,mark size=1.6pt] coordinates {(0,5/20)(6,10/20)(10,15/20)(21,20/20)};
\end{axis}\end{tikzpicture}\end{center}
\clearpage
\textbf{Question 3 --- Task 2 continued: comparison}\quad\hyperref[sec:q3]{Original description}
\par\smallskip\textbf{Same distribution: CDF above, PMF below.}
\begin{center}\begin{minipage}{\linewidth}\centering
\begin{tikzpicture}
\pgfplotsset{keySupport/.style={xmin=-2,xmax=23,xtick={0,6,10,21},
xlabel={$\GTenHLValue{x}$ (\GTenHLValue{USD})},
xticklabel={$\GTenHLValue{\pgfmathprintnumber{\tick}}$},
xticklabel style={font=\footnotesize,anchor=north,fill=white,inner sep=1pt},
xlabel style={fill=white,inner sep=2pt}}}
\begin{axis}[keyCombined,keySupport,name=keyTop,anchor=south west,
ylabel={$\HLCumulativeProbability{F_X(}\GTenHLValue{x}\HLCumulativeProbability{)}$},yticklabel={$\HLCumulativeProbability{\pgfmathprintnumber{\tick}}$}]
\keyPhantom{0}{5/20}{0/20}
\keyPhantom{6}{5/20}{5/20}
\keyPhantom{10}{5/20}{10/20}
\keyPhantom{21}{5/20}{15/20}
\addplot[keyStep,<-] coordinates {(-2,0)(0,0)};
\addplot[keyStep] coordinates {(0,5/20)(6,5/20)};
\addplot[keyStep] coordinates {(6,10/20)(10,10/20)};
\addplot[keyStep] coordinates {(10,15/20)(21,15/20)};
\addplot[keyStep,->] coordinates {(21,1)(23,1)};
\addplot[keyExcluded] coordinates {(0,0/20)(6,5/20)(10,10/20)(21,15/20)};
\addplot[keyIncluded] coordinates {(0,5/20)(6,10/20)(10,15/20)(21,20/20)};
\end{axis}
\begin{axis}[keyCombined,keySupport,at={(keyTop.south west)},
anchor=north west,yshift=-1.35cm,ylabel={$\GTenHLProbability{p_X(}\GTenHLValue{x}\GTenHLProbability{)}$},
yticklabel={$\GTenHLProbability{\pgfmathprintnumber{\tick}}$}]
\keyMass{0}{5/20}
\keyMass{6}{5/20}
\keyMass{10}{5/20}
\keyMass{21}{5/20}
\addlegendimage{keyStep}\addlegendentry{Top: \HLCumulativeProbability{CDF} steps}
\addlegendimage{keyPhantom,mark=*}\addlegendentry{Top: \GTenHLProbability{PMF} shadows (jumps)}
\addlegendimage{keyMass,mark=*}\addlegendentry{Bottom: \GTenHLProbability{PMF} stems}
\addlegendimage{keyCorrespondence}\addlegendentry{Midpoint correspondence}
\end{axis}
\keyCorrespondence{0}
\keyCorrespondence{6}
\keyCorrespondence{10}
\keyCorrespondence{21}
\end{tikzpicture}\end{minipage}\end{center}
\keyConnectionReading
\par\medskip
\textbf{CDF checks.}\begin{itemize}\item The intervals cover every real cutoff, and every height lies in $[0,1]$.\item The heights never decrease.\item Each segment includes its left endpoint and excludes its right endpoint; the filled point gives the value at a jump, so the CDF is right-continuous.\item The lower-tail limit is 0 and the upper-tail limit is 1; the arrows extend both tails.\item Successive cumulative differences recover the PMF numerators $5,5,5,5$ over $20$. Thus $\HLCumulativeProbability{F_X(}\GTenHLValue{a}\HLCumulativeProbability{)}-\HLCumulativeProbability{F_X(}\GTenHLValue{a^-}\HLCumulativeProbability{)}=\GTenHLProbability{p_X(}\GTenHLValue{a}\GTenHLProbability{)}$; unsupported cutoffs have zero jump and zero mass.\end{itemize}

\par\smallskip
\noindent\begin{minipage}{\linewidth}\centering
\textbf{Proposed partial credit.}\par\smallskip
\begin{tabular}{@{}l r@{}}\toprule
\CreditDescription{\textbf{Assessed step or criterion}} & \textbf{Points}\\\midrule
\CreditDescription{Cumulative sums} & 0.25\\
\CreditDescription{The complete interval table and tails} & 0.25\\
\CreditDescription{The labeled, correctly spaced step graph and endpoints} & 0.25\\
\CreditDescription{The CDF and jump checks} & 0.25\\
\midrule
\CreditDescription{\textbf{Total}} & \textbf{1}\\
\bottomrule\end{tabular}
\end{minipage}\par\smallskip

\par\vspace{3mm}\secInicio[][sec:q3]

\clearpage
\textbf{Question 3 --- continued}\quad\hyperref[sec:q3]{Original description}

\Needspace{8\baselineskip}
\subsection*{Task 3 (c) --- question and worked solution (proposed allocation) \score{0.5}}
\label{sol:q3t3}

\noindent\textbf{\boldmath Calculate and interpret (i) the probability that the handling fee is less than 10 USD and (ii) the probability that it is at least 21 USD.}\par\nopagebreak

\textbf{Allocation:} (i) 0.25 points; (ii) 0.25 points.

\textbf{(i)} Include only $0,6$ USD; exclude the boundary mass at 10 USD.
\[
\begin{aligned}
P(\GTenHLValue{X}<\GTenHLValue{10})&=\GTenHLProbability{p_X(}\GTenHLValue{0}\GTenHLProbability{)}+\GTenHLProbability{p_X(}\GTenHLValue{6}\GTenHLProbability{)}=\frac{5+5}{20}
\\ &=\GTenHLProbability{\frac{1}{2}}=\GTenHLProbability{0.50}=\GTenHLProbability{50\%}.
\end{aligned}
\]
The handling fee for one produce crate is below 10 USD with probability $\GTenHLProbability{1/2}$, or $\GTenHLProbability{50\%}$.

\textbf{(ii)} Include only $21$ USD; include the boundary mass at 21 USD.
\[
\begin{aligned}
P(\GTenHLValue{X}\geq\GTenHLValue{21})&=\GTenHLProbability{p_X(}\GTenHLValue{21}\GTenHLProbability{)}=\frac{5}{20}
\\ &=\GTenHLProbability{\frac{1}{4}}=\GTenHLProbability{0.25}=\GTenHLProbability{25\%}.
\end{aligned}
\]
The handling fee for one produce crate is at least 21 USD with probability $\GTenHLProbability{1/4}$, or $\GTenHLProbability{25\%}$.

\par\smallskip
\noindent\begin{minipage}{\linewidth}\centering
\textbf{Proposed partial credit.}\par\smallskip
\begin{tabular}{@{}l r@{}}\toprule
\CreditDescription{\textbf{Assessed step or criterion}} & \textbf{Points}\\\midrule
\CreditDescription{(i): selecting exactly the included PMF masses with the correct boundary} & 0.125\\
\CreditDescription{(i): summing them correctly} & 0.0625\\
\CreditDescription{(i): the contextual interpretation} & 0.0625\\
\CreditDescription{(ii): selecting exactly the included PMF masses with the correct boundary} & 0.125\\
\CreditDescription{(ii): summing them correctly} & 0.0625\\
\CreditDescription{(ii): the contextual interpretation} & 0.0625\\
\midrule
\CreditDescription{\textbf{Total}} & \textbf{0.5}\\
\bottomrule\end{tabular}
\end{minipage}\par\smallskip

Each part totals 0.25 points.

\par\vspace{3mm}\secInicio[][sec:q3]

\textbf{Question total: $0.5+1+0.5=2$ points for C2.}

\clearpage
\sectionScored{Equipment protection payouts}{C2-\sml{2}}{sec:q4}

An equipment protection provider classifies one claim in exactly one of eight exhaustive, mutually exclusive categories: no service, remote diagnosis, software reset, minor repair, major repair, replacement, store credit or cash refund. These eight named categories are the distinct outcomes in the sample space. Their probabilities are 0.15, 0.05, 0.15, 0.10, 0.10, 0.20, 0.10 and 0.15, respectively; do not assume equal likelihood. The corresponding payouts are 0, 0, 11, 11, 22, 22, 34 and 34 USD, respectively. Count store credit at its face value as a payout. Each payout is fixed for its category, and no additional payments apply. Record the provider's payout for one claim.

\Needspace{8\baselineskip}
\subsection*{Task 1 --- question and worked solution (proposed allocation) \score{0.5}}
\label{sol:q4t1}

\noindent\textbf{\boldmath Define the random variable and support; derive the grouped PMF; produce its table and a labeled graph with actual monetary spacing.}\par\nopagebreak

Let $\GTenHLValue{X}$ be the provider's payout for one claim, in USD. Map all eight categories as follows.
\begin{center}
\begin{tabular}{ccc}\toprule
Category & Given probability & $\GTenHLValue{X}$ (USD)\\\midrule
\GTenHLOutcome{No service} & $\GTenHLProbability{0.15}=\GTenHLProbability{3/20}$ & $\GTenHLValue{0}$\\
\GTenHLOutcome{Remote diagnosis} & $\GTenHLProbability{0.05}=\GTenHLProbability{1/20}$ & $\GTenHLValue{0}$\\
\GTenHLOutcome{Software reset} & $\GTenHLProbability{0.15}=\GTenHLProbability{3/20}$ & $\GTenHLValue{11}$\\
\GTenHLOutcome{Minor repair} & $\GTenHLProbability{0.10}=\GTenHLProbability{2/20}$ & $\GTenHLValue{11}$\\
\GTenHLOutcome{Major repair} & $\GTenHLProbability{0.10}=\GTenHLProbability{2/20}$ & $\GTenHLValue{22}$\\
\GTenHLOutcome{Replacement} & $\GTenHLProbability{0.20}=\GTenHLProbability{4/20}$ & $\GTenHLValue{22}$\\
\GTenHLOutcome{Store credit} & $\GTenHLProbability{0.10}=\GTenHLProbability{2/20}$ & $\GTenHLValue{34}$\\
\GTenHLOutcome{Cash refund} & $\GTenHLProbability{0.15}=\GTenHLProbability{3/20}$ & $\GTenHLValue{34}$\\
\bottomrule\end{tabular}
\end{center}
There are eight distinct sample-space outcomes and four monetary values. Add probabilities of mutually exclusive outcomes with the same value:

$\GTenHLValue{S_X}=\{0,11,22,34\}$ USD.
\par\vspace{3mm}\secInicio[][sec:q4]
\clearpage\textbf{Question 4 --- Task 1 continued: grouped \GTenHLProbability{PMF} and graph}
\begin{center}
\begingroup\setlength{\tabcolsep}{5pt}\renewcommand{\arraystretch}{1.18}
\begin{tabular}{@{}p{.30\linewidth}c p{.55\linewidth}@{}}\toprule
Complete outcomes or categories & $\GTenHLValue{x}$ (\GTenHLValue{USD}) & Grouped probability: $\GTenHLProbability{p_X(}\GTenHLValue{x}\GTenHLProbability{)}$\\\midrule
\GTenHLOutcome{No service}, \GTenHLOutcome{Remote diagnosis} & $\GTenHLValue{0}$ & $\begin{aligned}[t]\GTenHLProbability{p_X(}\GTenHLValue{0}\GTenHLProbability{)}&=\GTenHLProbability{3/20}+\GTenHLProbability{1/20}=\GTenHLProbability{4/20}\\&=\GTenHLProbability{0.2}=\GTenHLProbability{20\%}\end{aligned}$\\[5pt]
\GTenHLOutcome{Software reset}, \GTenHLOutcome{Minor repair} & $\GTenHLValue{11}$ & $\begin{aligned}[t]\GTenHLProbability{p_X(}\GTenHLValue{11}\GTenHLProbability{)}&=\GTenHLProbability{3/20}+\GTenHLProbability{2/20}=\GTenHLProbability{5/20}\\&=\GTenHLProbability{0.25}=\GTenHLProbability{25\%}\end{aligned}$\\[5pt]
\GTenHLOutcome{Major repair}, \GTenHLOutcome{Replacement} & $\GTenHLValue{22}$ & $\begin{aligned}[t]\GTenHLProbability{p_X(}\GTenHLValue{22}\GTenHLProbability{)}&=\GTenHLProbability{2/20}+\GTenHLProbability{4/20}=\GTenHLProbability{6/20}\\&=\GTenHLProbability{0.3}=\GTenHLProbability{30\%}\end{aligned}$\\[5pt]
\GTenHLOutcome{Store credit}, \GTenHLOutcome{Cash refund} & $\GTenHLValue{34}$ & $\begin{aligned}[t]\GTenHLProbability{p_X(}\GTenHLValue{34}\GTenHLProbability{)}&=\GTenHLProbability{2/20}+\GTenHLProbability{3/20}=\GTenHLProbability{5/20}\\&=\GTenHLProbability{0.25}=\GTenHLProbability{25\%}\end{aligned}$\\[5pt]
\bottomrule\end{tabular}
\endgroup
\end{center}
\begin{center}\begin{tikzpicture}
\begin{axis}[keyaxis,width=.86\linewidth,height=5cm,xmin=-2,xmax=36,ymin=0,ymax=.65,xtick={0,11,22,34},ylabel={$\GTenHLProbability{p_X(}\GTenHLValue{x}\GTenHLProbability{)}$},yticklabel={$\GTenHLProbability{\pgfmathprintnumber{\tick}}$}]
\draw[fill=keyMassFill,draw=keyMassColor] (axis cs:-0.45,0) rectangle (axis cs:0.45,4/20);
\draw[fill=keyMassFill,draw=keyMassColor] (axis cs:10.55,0) rectangle (axis cs:11.45,5/20);
\draw[fill=keyMassFill,draw=keyMassColor] (axis cs:21.55,0) rectangle (axis cs:22.45,6/20);
\draw[fill=keyMassFill,draw=keyMassColor] (axis cs:33.55,0) rectangle (axis cs:34.45,5/20);
\end{axis}\end{tikzpicture}\end{center}\Needspace{5\baselineskip}\textbf{PMF checks.}\begin{itemize}\item Every mass is positive.\item Normalization: $(4+5+6+5)/20=1$.\item Set $\GTenHLProbability{p_X(}\GTenHLValue{x}\GTenHLProbability{)}=\GTenHLProbability{0}$ outside the support. Bars use actual monetary spacing; gaps carry no mass.\end{itemize}

\textbf{Scoring note.} These explanatory positivity and normalization checks carry no separate points; omitting them does not reduce the task score.

\par\smallskip
\noindent\begin{minipage}{\linewidth}\centering
\textbf{Proposed partial credit.}\par\smallskip
\begin{tabular}{@{}l r@{}}\toprule
\CreditDescription{\textbf{Assessed step or criterion}} & \textbf{Points}\\\midrule
\CreditDescription{Variable, units and support} & 0.125\\
\CreditDescription{Category mapping, grouping and the PMF table} & 0.25\\
\CreditDescription{The labeled, correctly spaced graph} & 0.125\\
\midrule
\CreditDescription{\textbf{Total}} & \textbf{0.5}\\
\bottomrule\end{tabular}
\end{minipage}\par\smallskip

\par\vspace{3mm}\secInicio[][sec:q4]

\clearpage
\textbf{Question 4 --- continued}\quad\hyperref[sec:q4]{Original description}

\Needspace{8\baselineskip}
\subsection*{Task 2 --- question and worked solution (proposed allocation) \score{1}}
\label{sol:q4t2}

\noindent\textbf{\boldmath Construct the complete CDF table, including both tails, and draw a labeled step graph. Use open and filled endpoints correctly and preserve actual monetary spacing.}\par\nopagebreak

For $\HLCumulativeProbability{F_X(}\GTenHLValue{x}\HLCumulativeProbability{)}=P(X\leq x)$, add masses in increasing monetary order.

The successive cumulative numerators are $4,9,15,20$, all over $20$.
\begin{center}
\begingroup
\setlength{\tabcolsep}{5pt}
\begin{tabular}{@{}r@{}c@{}c@{}c@{}llc@{}}\toprule
\multicolumn{5}{c}{Cutoff interval (USD)} & Calculations & $\HLCumulativeProbability{F_X(}\GTenHLValue{x}\HLCumulativeProbability{)}$\\\midrule
 & & $\GTenHLValue{x}$ & $<$ & $\GTenHLValue{0}$ & $\text{No included support values}$ & $\HLCumulativeProbability{0}$\\[5pt]
$\GTenHLValue{0}$ & $\leq$ & $\GTenHLValue{x}$ & $<$ & $\GTenHLValue{11}$ & $\begin{aligned}[t]
  &\GTenHLProbability{p_X(}\GTenHLValue{0}\GTenHLProbability{)}
  \\&= \GTenHLProbability{\frac{4}{20}}
\end{aligned}$ & $\HLCumulativeProbability{4/20}$\\[5pt]
$\GTenHLValue{11}$ & $\leq$ & $\GTenHLValue{x}$ & $<$ & $\GTenHLValue{22}$ & $\begin{aligned}[t]
  &\GTenHLProbability{p_X(}\GTenHLValue{0}\GTenHLProbability{)} + \GTenHLProbability{p_X(}\GTenHLValue{11}\GTenHLProbability{)}
  \\&= \GTenHLProbability{\frac{4}{20}} + \GTenHLProbability{\frac{5}{20}}
\end{aligned}$ & $\HLCumulativeProbability{9/20}$\\[5pt]
$\GTenHLValue{22}$ & $\leq$ & $\GTenHLValue{x}$ & $<$ & $\GTenHLValue{34}$ & $\begin{aligned}[t]
  &\GTenHLProbability{p_X(}\GTenHLValue{0}\GTenHLProbability{)} + \GTenHLProbability{p_X(}\GTenHLValue{11}\GTenHLProbability{)} + \GTenHLProbability{p_X(}\GTenHLValue{22}\GTenHLProbability{)}
  \\&= \GTenHLProbability{\frac{4}{20}} + \GTenHLProbability{\frac{5}{20}} + \GTenHLProbability{\frac{6}{20}}
\end{aligned}$ & $\HLCumulativeProbability{15/20}$\\[5pt]
$\GTenHLValue{34}$ & $\leq$ & $\GTenHLValue{x}$ &  &  & $\begin{aligned}[t]
  &\GTenHLProbability{p_X(}\GTenHLValue{0}\GTenHLProbability{)} + \GTenHLProbability{p_X(}\GTenHLValue{11}\GTenHLProbability{)} + \GTenHLProbability{p_X(}\GTenHLValue{22}\GTenHLProbability{)} + \GTenHLProbability{p_X(}\GTenHLValue{34}\GTenHLProbability{)}
  \\&= \GTenHLProbability{\frac{4}{20}} + \GTenHLProbability{\frac{5}{20}} + \GTenHLProbability{\frac{6}{20}} + \GTenHLProbability{\frac{5}{20}}
\end{aligned}$ & $\HLCumulativeProbability{1}$\\[5pt]
\bottomrule\end{tabular}
\endgroup
\end{center}
\begin{center}\begin{tikzpicture}
\begin{axis}[keyaxis,width=.86\linewidth,height=7cm,xmin=-2,xmax=36,ymin=-.07,ymax=1.1,xtick={0,11,22,34},ylabel={$\HLCumulativeProbability{F_X(}\GTenHLValue{x}\HLCumulativeProbability{)}$},yticklabel={$\HLCumulativeProbability{\pgfmathprintnumber{\tick}}$},ytick={0,.25,.5,.75,1}]
\addplot[keyCDFColor,thick,<-] coordinates {(-2,0)(0,0)};
\addplot[keyCDFColor,thick] coordinates {(0,4/20)(11,4/20)};
\addplot[keyCDFColor,thick] coordinates {(11,9/20)(22,9/20)};
\addplot[keyCDFColor,thick] coordinates {(22,15/20)(34,15/20)};
\addplot[keyCDFColor,thick,->] coordinates {(34,1)(36,1)};
\addplot[keyCDFColor,only marks,mark=*,mark size=1.6pt,mark options={fill=white}] coordinates {(0,0/20)(11,4/20)(22,9/20)(34,15/20)};
\addplot[keyCDFColor,only marks,mark=*,mark size=1.6pt] coordinates {(0,4/20)(11,9/20)(22,15/20)(34,20/20)};
\end{axis}\end{tikzpicture}\end{center}
\clearpage
\textbf{Question 4 --- Task 2 continued: comparison}\quad\hyperref[sec:q4]{Original description}
\par\smallskip\textbf{Same distribution: CDF above, PMF below.}
\begin{center}\begin{minipage}{\linewidth}\centering
\begin{tikzpicture}
\pgfplotsset{keySupport/.style={xmin=-2,xmax=36,xtick={0,11,22,34},
xlabel={$\GTenHLValue{x}$ (\GTenHLValue{USD})},
xticklabel={$\GTenHLValue{\pgfmathprintnumber{\tick}}$},
xticklabel style={font=\footnotesize,anchor=north,fill=white,inner sep=1pt},
xlabel style={fill=white,inner sep=2pt}}}
\begin{axis}[keyCombined,keySupport,name=keyTop,anchor=south west,
ylabel={$\HLCumulativeProbability{F_X(}\GTenHLValue{x}\HLCumulativeProbability{)}$},yticklabel={$\HLCumulativeProbability{\pgfmathprintnumber{\tick}}$}]
\keyPhantom{0}{4/20}{0/20}
\keyPhantom{11}{5/20}{4/20}
\keyPhantom{22}{6/20}{9/20}
\keyPhantom{34}{5/20}{15/20}
\addplot[keyStep,<-] coordinates {(-2,0)(0,0)};
\addplot[keyStep] coordinates {(0,4/20)(11,4/20)};
\addplot[keyStep] coordinates {(11,9/20)(22,9/20)};
\addplot[keyStep] coordinates {(22,15/20)(34,15/20)};
\addplot[keyStep,->] coordinates {(34,1)(36,1)};
\addplot[keyExcluded] coordinates {(0,0/20)(11,4/20)(22,9/20)(34,15/20)};
\addplot[keyIncluded] coordinates {(0,4/20)(11,9/20)(22,15/20)(34,20/20)};
\end{axis}
\begin{axis}[keyCombined,keySupport,at={(keyTop.south west)},
anchor=north west,yshift=-1.35cm,ylabel={$\GTenHLProbability{p_X(}\GTenHLValue{x}\GTenHLProbability{)}$},
yticklabel={$\GTenHLProbability{\pgfmathprintnumber{\tick}}$}]
\keyMass{0}{4/20}
\keyMass{11}{5/20}
\keyMass{22}{6/20}
\keyMass{34}{5/20}
\addlegendimage{keyStep}\addlegendentry{Top: \HLCumulativeProbability{CDF} steps}
\addlegendimage{keyPhantom,mark=*}\addlegendentry{Top: \GTenHLProbability{PMF} shadows (jumps)}
\addlegendimage{keyMass,mark=*}\addlegendentry{Bottom: \GTenHLProbability{PMF} stems}
\addlegendimage{keyCorrespondence}\addlegendentry{Midpoint correspondence}
\end{axis}
\keyCorrespondence{0}
\keyCorrespondence{11}
\keyCorrespondence{22}
\keyCorrespondence{34}
\end{tikzpicture}\end{minipage}\end{center}
\keyConnectionReading
\par\medskip
\textbf{CDF checks.}\begin{itemize}\item The intervals cover every real cutoff, and every height lies in $[0,1]$.\item The heights never decrease.\item Each segment includes its left endpoint and excludes its right endpoint; the filled point gives the value at a jump, so the CDF is right-continuous.\item The lower-tail limit is 0 and the upper-tail limit is 1; the arrows extend both tails.\item Successive cumulative differences recover the PMF numerators $4,5,6,5$ over $20$. Thus $\HLCumulativeProbability{F_X(}\GTenHLValue{a}\HLCumulativeProbability{)}-\HLCumulativeProbability{F_X(}\GTenHLValue{a^-}\HLCumulativeProbability{)}=\GTenHLProbability{p_X(}\GTenHLValue{a}\GTenHLProbability{)}$; unsupported cutoffs have zero jump and zero mass.\end{itemize}

\par\smallskip
\noindent\begin{minipage}{\linewidth}\centering
\textbf{Proposed partial credit.}\par\smallskip
\begin{tabular}{@{}l r@{}}\toprule
\CreditDescription{\textbf{Assessed step or criterion}} & \textbf{Points}\\\midrule
\CreditDescription{Cumulative sums} & 0.25\\
\CreditDescription{The complete interval table and tails} & 0.25\\
\CreditDescription{The labeled, correctly spaced step graph and endpoints} & 0.25\\
\CreditDescription{The CDF and jump checks} & 0.25\\
\midrule
\CreditDescription{\textbf{Total}} & \textbf{1}\\
\bottomrule\end{tabular}
\end{minipage}\par\smallskip

\par\vspace{3mm}\secInicio[][sec:q4]

\clearpage
\textbf{Question 4 --- continued}\quad\hyperref[sec:q4]{Original description}

\Needspace{8\baselineskip}
\subsection*{Task 3 (c) --- question and worked solution (proposed allocation) \score{0.5}}
\label{sol:q4t3}

\noindent\textbf{\boldmath Calculate and interpret (i) the probability that the provider's payout is less than 11 USD and (ii) the probability that it is at least 22 USD.}\par\nopagebreak

\textbf{Allocation:} (i) 0.25 points; (ii) 0.25 points.

\textbf{(i)} Include only $0$ USD; exclude the boundary mass at 11 USD.
\[
\begin{aligned}
P(\GTenHLValue{X}<\GTenHLValue{11})&=\GTenHLProbability{p_X(}\GTenHLValue{0}\GTenHLProbability{)}=\frac{4}{20}
\\ &=\GTenHLProbability{\frac{1}{5}}=\GTenHLProbability{0.20}=\GTenHLProbability{20\%}.
\end{aligned}
\]
The provider's payout for one claim is below 11 USD with probability $\GTenHLProbability{1/5}$, or $\GTenHLProbability{20\%}$.

\textbf{(ii)} Include only $22,34$ USD; include the boundary mass at 22 USD.
\[
\begin{aligned}
P(\GTenHLValue{X}\geq\GTenHLValue{22})&=\GTenHLProbability{p_X(}\GTenHLValue{22}\GTenHLProbability{)}+\GTenHLProbability{p_X(}\GTenHLValue{34}\GTenHLProbability{)}=\frac{6+5}{20}
\\ &=\GTenHLProbability{\frac{11}{20}}=\GTenHLProbability{0.55}=\GTenHLProbability{55\%}.
\end{aligned}
\]
The provider's payout for one claim is at least 22 USD with probability $\GTenHLProbability{11/20}$, or $\GTenHLProbability{55\%}$.

\par\smallskip
\noindent\begin{minipage}{\linewidth}\centering
\textbf{Proposed partial credit.}\par\smallskip
\begin{tabular}{@{}l r@{}}\toprule
\CreditDescription{\textbf{Assessed step or criterion}} & \textbf{Points}\\\midrule
\CreditDescription{(i): selecting exactly the included PMF masses with the correct boundary} & 0.125\\
\CreditDescription{(i): summing them correctly} & 0.0625\\
\CreditDescription{(i): the contextual interpretation} & 0.0625\\
\CreditDescription{(ii): selecting exactly the included PMF masses with the correct boundary} & 0.125\\
\CreditDescription{(ii): summing them correctly} & 0.0625\\
\CreditDescription{(ii): the contextual interpretation} & 0.0625\\
\midrule
\CreditDescription{\textbf{Total}} & \textbf{0.5}\\
\bottomrule\end{tabular}
\end{minipage}\par\smallskip

Each part totals 0.25 points.

\par\vspace{3mm}\secInicio[][sec:q4]

\textbf{Question total: $0.5+1+0.5=2$ points for C2.}

\clearpage
\sectionScored{Travel platform account credits}{C2-\sml{4}}{sec:q5}

For one booking, a travel platform records exactly one of seven distinct, mutually exclusive and exhaustive booking-resolution outcomes: A (booking completed as scheduled); B (customer cancellation within the free period); C (customer cancellation after the free period); D (duplicate booking corrected); E (alternative accommodation accepted); F (rescheduled journey accepted); G (provider cancellation with refund). These named outcomes form the sample space. Outcomes A and B award 0 USD account credit; C and D award 12 USD; E, F and G award 18 USD. Thus seven outcomes map to three possible credit values. B is twice as likely as A; C is three times as likely as A; D is twice as likely as B; E is as likely as B; F is twice as likely as E; G is as likely as F. There are no other payments or adjustments. Determine the outcome probabilities from these likelihood relationships and record the account credit.

\Needspace{8\baselineskip}
\subsection*{Task 1 --- question and worked solution (proposed allocation) \score{1}}
\label{sol:q5t1}

\noindent\textbf{\boldmath Define the random variable and support. Derive the seven named outcome probabilities using normalization, map each outcome to its credit value, and construct the grouped PMF table and labeled graph.}\par\nopagebreak

Let $\GTenHLValue{X}$ be the account credit for one booking, in USD. The sample space is $\Omega=\{A,B,C,D,E,F,G\}$, whereas $\GTenHLValue{S_X}=\{0,12,18\}$ USD. Set $\GTenHLParameter{a}=P(A)$, an outcome probability, not the grouped mass $\GTenHLProbability{p_X(}\GTenHLValue{0}\GTenHLProbability{)}$.
\begin{align*}
P(A)&=a,&P(B)&=2a,&P(C)&=3a,\\
P(D)&=2P(B)=4a,&P(E)&=P(B)=2a,\\
P(F)&=2P(E)=4a,&P(G)&=P(F)=4a.
\end{align*}
Normalization over the seven exhaustive outcomes gives
\[
a+2a+3a+4a+2a+4a+4a=1,\qquad 20a=1,\qquad \GTenHLParameter{a}=\frac1{20}.
\]
The complete outcome probabilities and mappings are:
\begin{center}
\begin{tabular}{ccc}\toprule
Named outcome & $\GTenHLProbability{P(}\GTenHLOutcome{\mathrm{outcome}}\GTenHLProbability{)}$ & $\GTenHLValue{X}$ (USD)\\
\midrule
\GTenHLOutcome{A: Booking completed as scheduled} & $\GTenHLParameter{a}=\GTenHLProbability{1/20}$ & $\GTenHLValue{0}$\\
\GTenHLOutcome{B: Customer cancellation within the free period} & $2\GTenHLParameter{a}=\GTenHLProbability{2/20}$ & $\GTenHLValue{0}$\\
\GTenHLOutcome{C: Customer cancellation after the free period} & $3\GTenHLParameter{a}=\GTenHLProbability{3/20}$ & $\GTenHLValue{12}$\\
\GTenHLOutcome{D: Duplicate booking corrected} & $4\GTenHLParameter{a}=\GTenHLProbability{4/20}$ & $\GTenHLValue{12}$\\
\GTenHLOutcome{E: Alternative accommodation accepted} & $2\GTenHLParameter{a}=\GTenHLProbability{2/20}$ & $\GTenHLValue{18}$\\
\GTenHLOutcome{F: Rescheduled journey accepted} & $4\GTenHLParameter{a}=\GTenHLProbability{4/20}$ & $\GTenHLValue{18}$\\
\GTenHLOutcome{G: Provider cancellation with refund} & $4\GTenHLParameter{a}=\GTenHLProbability{4/20}$ & $\GTenHLValue{18}$\\
\bottomrule\end{tabular}
\end{center}
\par\vspace{3mm}\secInicio[][sec:q5]
\clearpage\textbf{Question 5 --- Task 1 continued: outcome aggregation and grouped PMF}

The seven outcomes are distinct even when they award the same credit. Group their probabilities to obtain $\GTenHLProbability{p_X(}\GTenHLValue{x}\GTenHLProbability{)}=P(X=x)$:

$\GTenHLValue{S_X}=\{0,12,18\}$ USD.
\begin{center}
\begingroup\setlength{\tabcolsep}{5pt}\renewcommand{\arraystretch}{1.18}
\begin{tabular}{@{}p{.30\linewidth}c p{.55\linewidth}@{}}\toprule
Complete outcomes or categories & $\GTenHLValue{x}$ (\GTenHLValue{USD}) & Grouped probability: $\GTenHLProbability{p_X(}\GTenHLValue{x}\GTenHLProbability{)}$\\\midrule
\GTenHLOutcome{A}, \GTenHLOutcome{B} & $\GTenHLValue{0}$ & $\begin{aligned}[t]\GTenHLProbability{p_X(}\GTenHLValue{0}\GTenHLProbability{)}&=\GTenHLProbability{1/20}+\GTenHLProbability{2/20}=\GTenHLProbability{3/20}\\&=\GTenHLProbability{0.15}=\GTenHLProbability{15\%}\end{aligned}$\\[5pt]
\GTenHLOutcome{C}, \GTenHLOutcome{D} & $\GTenHLValue{12}$ & $\begin{aligned}[t]\GTenHLProbability{p_X(}\GTenHLValue{12}\GTenHLProbability{)}&=\GTenHLProbability{3/20}+\GTenHLProbability{4/20}=\GTenHLProbability{7/20}\\&=\GTenHLProbability{0.35}=\GTenHLProbability{35\%}\end{aligned}$\\[5pt]
\GTenHLOutcome{E}, \GTenHLOutcome{F}, \GTenHLOutcome{G} & $\GTenHLValue{18}$ & $\begin{aligned}[t]\GTenHLProbability{p_X(}\GTenHLValue{18}\GTenHLProbability{)}&=\GTenHLProbability{2/20}+\GTenHLProbability{4/20}+\GTenHLProbability{4/20}=\GTenHLProbability{10/20}\\&=\GTenHLProbability{0.5}=\GTenHLProbability{50\%}\end{aligned}$\\[5pt]
\bottomrule\end{tabular}
\endgroup
\end{center}
\begin{center}\begin{tikzpicture}
\begin{axis}[keyaxis,width=.86\linewidth,height=5cm,xmin=-2,xmax=20,ymin=0,ymax=.65,xtick={0,12,18},ylabel={$\GTenHLProbability{p_X(}\GTenHLValue{x}\GTenHLProbability{)}$},yticklabel={$\GTenHLProbability{\pgfmathprintnumber{\tick}}$}]
\draw[fill=keyMassFill,draw=keyMassColor] (axis cs:-0.45,0) rectangle (axis cs:0.45,3/20);
\draw[fill=keyMassFill,draw=keyMassColor] (axis cs:11.55,0) rectangle (axis cs:12.45,7/20);
\draw[fill=keyMassFill,draw=keyMassColor] (axis cs:17.55,0) rectangle (axis cs:18.45,10/20);
\end{axis}\end{tikzpicture}\end{center}\Needspace{5\baselineskip}\textbf{PMF checks.}\begin{itemize}\item Every mass is positive.\item Normalization: $(3+7+10)/20=1$.\item Set $\GTenHLProbability{p_X(}\GTenHLValue{x}\GTenHLProbability{)}=\GTenHLProbability{0}$ outside the support. Bars use actual monetary spacing; gaps carry no mass.\end{itemize}
\textbf{Relationship check.} Outcome numerators satisfy $2=2(1)$, $3=3(1)$, $4=2(2)$, $2=2$, $4=2(2)$ and $4=4$, reproducing all six relationships.

\textbf{Scoring note.} These explanatory positivity and normalization checks carry no separate points; omitting them does not reduce the task score.

\par\smallskip
\noindent\begin{minipage}{\linewidth}\centering
\textbf{Proposed partial credit.}\par\smallskip
\begin{tabular}{@{}l r@{}}\toprule
\CreditDescription{\textbf{Assessed step or criterion}} & \textbf{Points}\\\midrule
\CreditDescription{Variable, units and support} & 0.125\\
\CreditDescription{Translating the relationships and solving normalization} & 0.375\\
\CreditDescription{All seven outcome probabilities, mappings and three grouped PMF entries} & 0.25\\
\CreditDescription{The labeled, correctly spaced graph} & 0.125\\
\CreditDescription{Relationship checks} & 0.125\\
\midrule
\CreditDescription{\textbf{Total}} & \textbf{1}\\
\bottomrule\end{tabular}
\end{minipage}\par\smallskip

\par\vspace{3mm}\secInicio[][sec:q5]

\clearpage
\textbf{Question 5 --- continued}\quad\hyperref[sec:q5]{Original description}

\Needspace{8\baselineskip}
\subsection*{Task 2 --- question and worked solution (proposed allocation) \score{2}}
\label{sol:q5t2}

\noindent\textbf{\boldmath Construct the complete CDF table, including both tails, and draw a labeled step graph. Use open and filled endpoints correctly and preserve actual monetary spacing.}\par\nopagebreak

For $\HLCumulativeProbability{F_X(}\GTenHLValue{x}\HLCumulativeProbability{)}=P(X\leq x)$, add masses in increasing monetary order.

The successive cumulative numerators are $3,10,20$, all over $20$.
\begin{center}
\begingroup
\setlength{\tabcolsep}{5pt}
\begin{tabular}{@{}r@{}c@{}c@{}c@{}llc@{}}\toprule
\multicolumn{5}{c}{Cutoff interval (USD)} & Calculations & $\HLCumulativeProbability{F_X(}\GTenHLValue{x}\HLCumulativeProbability{)}$\\\midrule
 & & $\GTenHLValue{x}$ & $<$ & $\GTenHLValue{0}$ & $\text{No included support values}$ & $\HLCumulativeProbability{0}$\\[5pt]
$\GTenHLValue{0}$ & $\leq$ & $\GTenHLValue{x}$ & $<$ & $\GTenHLValue{12}$ & $\begin{aligned}[t]
  &\GTenHLProbability{p_X(}\GTenHLValue{0}\GTenHLProbability{)}
  \\&= \GTenHLProbability{\frac{3}{20}}
\end{aligned}$ & $\HLCumulativeProbability{3/20}$\\[5pt]
$\GTenHLValue{12}$ & $\leq$ & $\GTenHLValue{x}$ & $<$ & $\GTenHLValue{18}$ & $\begin{aligned}[t]
  &\GTenHLProbability{p_X(}\GTenHLValue{0}\GTenHLProbability{)} + \GTenHLProbability{p_X(}\GTenHLValue{12}\GTenHLProbability{)}
  \\&= \GTenHLProbability{\frac{3}{20}} + \GTenHLProbability{\frac{7}{20}}
\end{aligned}$ & $\HLCumulativeProbability{10/20}$\\[5pt]
$\GTenHLValue{18}$ & $\leq$ & $\GTenHLValue{x}$ &  &  & $\begin{aligned}[t]
  &\GTenHLProbability{p_X(}\GTenHLValue{0}\GTenHLProbability{)} + \GTenHLProbability{p_X(}\GTenHLValue{12}\GTenHLProbability{)} + \GTenHLProbability{p_X(}\GTenHLValue{18}\GTenHLProbability{)}
  \\&= \GTenHLProbability{\frac{3}{20}} + \GTenHLProbability{\frac{7}{20}} + \GTenHLProbability{\frac{10}{20}}
\end{aligned}$ & $\HLCumulativeProbability{1}$\\[5pt]
\bottomrule\end{tabular}
\endgroup
\end{center}
\begin{center}\begin{tikzpicture}
\begin{axis}[keyaxis,width=.86\linewidth,height=7cm,xmin=-2,xmax=20,ymin=-.07,ymax=1.1,xtick={0,12,18},ylabel={$\HLCumulativeProbability{F_X(}\GTenHLValue{x}\HLCumulativeProbability{)}$},yticklabel={$\HLCumulativeProbability{\pgfmathprintnumber{\tick}}$},ytick={0,.25,.5,.75,1}]
\addplot[keyCDFColor,thick,<-] coordinates {(-2,0)(0,0)};
\addplot[keyCDFColor,thick] coordinates {(0,3/20)(12,3/20)};
\addplot[keyCDFColor,thick] coordinates {(12,10/20)(18,10/20)};
\addplot[keyCDFColor,thick,->] coordinates {(18,1)(20,1)};
\addplot[keyCDFColor,only marks,mark=*,mark size=1.6pt,mark options={fill=white}] coordinates {(0,0/20)(12,3/20)(18,10/20)};
\addplot[keyCDFColor,only marks,mark=*,mark size=1.6pt] coordinates {(0,3/20)(12,10/20)(18,20/20)};
\end{axis}\end{tikzpicture}\end{center}
\clearpage
\textbf{Question 5 --- Task 2 continued: comparison}\quad\hyperref[sec:q5]{Original description}
\par\smallskip\textbf{Same distribution: CDF above, PMF below.}
\begin{center}\begin{minipage}{\linewidth}\centering
\begin{tikzpicture}
\pgfplotsset{keySupport/.style={xmin=-2,xmax=20,xtick={0,12,18},
xlabel={$\GTenHLValue{x}$ (\GTenHLValue{USD})},
xticklabel={$\GTenHLValue{\pgfmathprintnumber{\tick}}$},
xticklabel style={font=\footnotesize,anchor=north,fill=white,inner sep=1pt},
xlabel style={fill=white,inner sep=2pt}}}
\begin{axis}[keyCombined,keySupport,name=keyTop,anchor=south west,
ylabel={$\HLCumulativeProbability{F_X(}\GTenHLValue{x}\HLCumulativeProbability{)}$},yticklabel={$\HLCumulativeProbability{\pgfmathprintnumber{\tick}}$}]
\keyPhantom{0}{3/20}{0/20}
\keyPhantom{12}{7/20}{3/20}
\keyPhantom{18}{10/20}{10/20}
\addplot[keyStep,<-] coordinates {(-2,0)(0,0)};
\addplot[keyStep] coordinates {(0,3/20)(12,3/20)};
\addplot[keyStep] coordinates {(12,10/20)(18,10/20)};
\addplot[keyStep,->] coordinates {(18,1)(20,1)};
\addplot[keyExcluded] coordinates {(0,0/20)(12,3/20)(18,10/20)};
\addplot[keyIncluded] coordinates {(0,3/20)(12,10/20)(18,20/20)};
\end{axis}
\begin{axis}[keyCombined,keySupport,at={(keyTop.south west)},
anchor=north west,yshift=-1.35cm,ylabel={$\GTenHLProbability{p_X(}\GTenHLValue{x}\GTenHLProbability{)}$},
yticklabel={$\GTenHLProbability{\pgfmathprintnumber{\tick}}$}]
\keyMass{0}{3/20}
\keyMass{12}{7/20}
\keyMass{18}{10/20}
\addlegendimage{keyStep}\addlegendentry{Top: \HLCumulativeProbability{CDF} steps}
\addlegendimage{keyPhantom,mark=*}\addlegendentry{Top: \GTenHLProbability{PMF} shadows (jumps)}
\addlegendimage{keyMass,mark=*}\addlegendentry{Bottom: \GTenHLProbability{PMF} stems}
\addlegendimage{keyCorrespondence}\addlegendentry{Midpoint correspondence}
\end{axis}
\keyCorrespondence{0}
\keyCorrespondence{12}
\keyCorrespondence{18}
\end{tikzpicture}\end{minipage}\end{center}
\keyConnectionReading
\par\medskip
\textbf{CDF checks.}\begin{itemize}\item The intervals cover every real cutoff, and every height lies in $[0,1]$.\item The heights never decrease.\item Each segment includes its left endpoint and excludes its right endpoint; the filled point gives the value at a jump, so the CDF is right-continuous.\item The lower-tail limit is 0 and the upper-tail limit is 1; the arrows extend both tails.\item Successive cumulative differences recover the PMF numerators $3,7,10$ over $20$. Thus $\HLCumulativeProbability{F_X(}\GTenHLValue{a}\HLCumulativeProbability{)}-\HLCumulativeProbability{F_X(}\GTenHLValue{a^-}\HLCumulativeProbability{)}=\GTenHLProbability{p_X(}\GTenHLValue{a}\GTenHLProbability{)}$; unsupported cutoffs have zero jump and zero mass.\end{itemize}

\par\smallskip
\noindent\begin{minipage}{\linewidth}\centering
\textbf{Proposed partial credit.}\par\smallskip
\begin{tabular}{@{}l r@{}}\toprule
\CreditDescription{\textbf{Assessed step or criterion}} & \textbf{Points}\\\midrule
\CreditDescription{Cumulative sums} & 0.5\\
\CreditDescription{The complete interval table and tails} & 0.5\\
\CreditDescription{The labeled, correctly spaced step graph and endpoints} & 0.5\\
\CreditDescription{The CDF and jump checks} & 0.5\\
\midrule
\CreditDescription{\textbf{Total}} & \textbf{2}\\
\bottomrule\end{tabular}
\end{minipage}\par\smallskip

\par\vspace{3mm}\secInicio[][sec:q5]

\clearpage
\textbf{Question 5 --- continued}\quad\hyperref[sec:q5]{Original description}

\Needspace{8\baselineskip}
\subsection*{Task 3 (c) --- question and worked solution (proposed allocation) \score{1}}
\label{sol:q5t3}

\noindent\textbf{\boldmath Calculate and interpret (i) the probability that the account credit is less than 12 USD and (ii) the probability that it is at least 18 USD.}\par\nopagebreak

\textbf{Allocation:} (i) 0.5 points; (ii) 0.5 points.

\textbf{(i)} Include only $0$ USD; exclude the boundary mass at 12 USD.
\[
\begin{aligned}
P(\GTenHLValue{X}<\GTenHLValue{12})&=\GTenHLProbability{p_X(}\GTenHLValue{0}\GTenHLProbability{)}=\frac{3}{20}
\\ &=\GTenHLProbability{\frac{3}{20}}=\GTenHLProbability{0.15}=\GTenHLProbability{15\%}.
\end{aligned}
\]
The account credit for one booking is below 12 USD with probability $\GTenHLProbability{3/20}$, or $\GTenHLProbability{15\%}$.

\textbf{(ii)} Include only $18$ USD; include the boundary mass at 18 USD.
\[
\begin{aligned}
P(\GTenHLValue{X}\geq\GTenHLValue{18})&=\GTenHLProbability{p_X(}\GTenHLValue{18}\GTenHLProbability{)}=\frac{10}{20}
\\ &=\GTenHLProbability{\frac{1}{2}}=\GTenHLProbability{0.50}=\GTenHLProbability{50\%}.
\end{aligned}
\]
The account credit for one booking is at least 18 USD with probability $\GTenHLProbability{1/2}$, or $\GTenHLProbability{50\%}$.

\par\smallskip
\noindent\begin{minipage}{\linewidth}\centering
\textbf{Proposed partial credit.}\par\smallskip
\begin{tabular}{@{}l r@{}}\toprule
\CreditDescription{\textbf{Assessed step or criterion}} & \textbf{Points}\\\midrule
\CreditDescription{(i): selecting exactly the included PMF masses with the correct boundary} & 0.25\\
\CreditDescription{(i): summing them correctly} & 0.125\\
\CreditDescription{(i): the contextual interpretation} & 0.125\\
\CreditDescription{(ii): selecting exactly the included PMF masses with the correct boundary} & 0.25\\
\CreditDescription{(ii): summing them correctly} & 0.125\\
\CreditDescription{(ii): the contextual interpretation} & 0.125\\
\midrule
\CreditDescription{\textbf{Total}} & \textbf{1}\\
\bottomrule\end{tabular}
\end{minipage}\par\smallskip

Each part totals 0.5 points.

\par\vspace{3mm}\secInicio[][sec:q5]

\textbf{Question total: $1+2+1=4$ points for C2.}

\clearpage
\sectionScored{Library fundraiser prize tickets}{C2-\sml{6}}{sec:q6}

A library fundraiser has nine physically distinct prize tickets: three worth 0 USD each, four worth 6 USD each and two worth 14 USD each. Choose a subset of three tickets uniformly without replacement, so every subset of three physically distinct tickets has the same chance. Ignore selection order. First add the face values. Then add one bonus of 8 USD if the selected subset contains at least one 6 USD ticket and at least one 14 USD ticket; otherwise add no bonus. Finally, if this adjusted amount exceeds 30 USD, pay only 30 USD; otherwise pay the adjusted amount. No other payments or deductions apply. Record the final payout.

\Needspace{8\baselineskip}
\subsection*{Task 1 --- question and worked solution (proposed allocation) \score{1.5}}
\label{sol:q6t1}

\noindent\textbf{\boldmath Define the random variable and support. Count the distinct-ticket selections, apply the bonus and then the cap, aggregate equal final payouts, and derive the PMF. Produce its table and labeled graph.}\par\nopagebreak

Let $\GTenHLValue{X}$ be the final payout for one three-ticket subset, in USD. Label the distinct tickets $Z_1,Z_2,Z_3$, $L_1,\ldots,L_4$, and $H_1,H_2$, worth 0, 6 and 14 USD respectively. There are $\binom93=84$ equally likely unordered subsets, each with probability $1/84$.

A composition $(a,b,c)$ counts zero-, six- and fourteen-dollar tickets. It satisfies $a+b+c=3$, $0\leq a\leq3$, $0\leq b\leq4$ and $0\leq c\leq2$. Its subset count is $N(a,b,c)=\binom3a\binom4b\binom2c$ and its probability is $N(a,b,c)/84$. Let $I=1$ when $b\geq1$ and $c\geq1$, and 0 otherwise. First calculate $Y=6b+14c+8I$, then $X=\min\{Y,30\}$: the bonus is added once, before the cap.
\begingroup\renewcommand{\arraystretch}{1.35}

\begin{center}
\begin{tabular}{ccccccc}\toprule
$(a,b,c)$ & Subset count & Probability & Face sum & Bonus & $Y$ & $X$\\\midrule
$(\GTenHLOutcome{3},\GTenHLOutcome{0},\GTenHLOutcome{0})$ & $\binom3{3}\binom4{0}\binom2{0}=1$ & $\GTenHLProbability{1/84}$ & $\GTenHLValue{0}$ & $\GTenHLValue{0}$ & $\GTenHLSecond{0}$ & $\GTenHLValue{0}$\\
$(\GTenHLOutcome{2},\GTenHLOutcome{1},\GTenHLOutcome{0})$ & $\binom3{2}\binom4{1}\binom2{0}=12$ & $\GTenHLProbability{12/84}$ & $\GTenHLValue{6}$ & $\GTenHLValue{0}$ & $\GTenHLSecond{6}$ & $\GTenHLValue{6}$\\
$(\GTenHLOutcome{1},\GTenHLOutcome{2},\GTenHLOutcome{0})$ & $\binom3{1}\binom4{2}\binom2{0}=18$ & $\GTenHLProbability{18/84}$ & $\GTenHLValue{12}$ & $\GTenHLValue{0}$ & $\GTenHLSecond{12}$ & $\GTenHLValue{12}$\\
$(\GTenHLOutcome{0},\GTenHLOutcome{3},\GTenHLOutcome{0})$ & $\binom3{0}\binom4{3}\binom2{0}=4$ & $\GTenHLProbability{4/84}$ & $\GTenHLValue{18}$ & $\GTenHLValue{0}$ & $\GTenHLSecond{18}$ & $\GTenHLValue{18}$\\
$(\GTenHLOutcome{2},\GTenHLOutcome{0},\GTenHLOutcome{1})$ & $\binom3{2}\binom4{0}\binom2{1}=6$ & $\GTenHLProbability{6/84}$ & $\GTenHLValue{14}$ & $\GTenHLValue{0}$ & $\GTenHLSecond{14}$ & $\GTenHLValue{14}$\\
$(\GTenHLOutcome{1},\GTenHLOutcome{1},\GTenHLOutcome{1})$ & $\binom3{1}\binom4{1}\binom2{1}=24$ & $\GTenHLProbability{24/84}$ & $\GTenHLValue{20}$ & $\GTenHLValue{8}$ & $\GTenHLSecond{28}$ & $\GTenHLValue{28}$\\
$(\GTenHLOutcome{0},\GTenHLOutcome{2},\GTenHLOutcome{1})$ & $\binom3{0}\binom4{2}\binom2{1}=12$ & $\GTenHLProbability{12/84}$ & $\GTenHLValue{26}$ & $\GTenHLValue{8}$ & $\GTenHLSecond{34}$ & $\GTenHLValue{30}$\\
$(\GTenHLOutcome{1},\GTenHLOutcome{0},\GTenHLOutcome{2})$ & $\binom3{1}\binom4{0}\binom2{2}=3$ & $\GTenHLProbability{3/84}$ & $\GTenHLValue{28}$ & $\GTenHLValue{0}$ & $\GTenHLSecond{28}$ & $\GTenHLValue{28}$\\
$(\GTenHLOutcome{0},\GTenHLOutcome{1},\GTenHLOutcome{2})$ & $\binom3{0}\binom4{1}\binom2{2}=4$ & $\GTenHLProbability{4/84}$ & $\GTenHLValue{34}$ & $\GTenHLValue{8}$ & $\GTenHLSecond{42}$ & $\GTenHLValue{30}$\\
\bottomrule\end{tabular}
\end{center}
\endgroup
All monetary columns are in USD. Listing $c=0,1,2$ and each feasible $b$ gives all nine compositions. Their counts sum to $1+12+18+4+6+24+12+3+4=84$. Compositions are not equally likely: $(1,1,1)$ represents 24 subsets, while $(3,0,0)$ represents only one.

\textbf{Aggregation.} Below the cap, the 28 USD payout combines $(1,1,1)$ (24 subsets after the bonus) with $(1,0,2)$ (3 subsets without a bonus), giving $\GTenHLProbability{p_X(}\GTenHLValue{28}\GTenHLProbability{)}=(24+3)/84=27/84$. At the cap, adjusted amounts 34 and 42 USD from $(0,2,1)$ and $(0,1,2)$ combine, giving $\GTenHLProbability{p_X(}\GTenHLValue{30}\GTenHLProbability{)}=(12+4)/84=16/84$. Every other payout comes from one composition. Thus $S_X=\{0,6,12,14,18,28,30\}$ USD; the 28 USD payout stays below the cap.
\par\vspace{3mm}\secInicio[][sec:q6]
\clearpage\textbf{Question 6 --- Task 1 continued: grouped \GTenHLProbability{PMF} and graph}

Each composition $(a,b,c)$ represents the stated number of distinct-ticket subsets; the original composition table supplies its binomial count. Add their probabilities when final payouts agree.
\begin{center}
\begingroup\setlength{\tabcolsep}{5pt}\renewcommand{\arraystretch}{1.18}
\begin{tabular}{@{}p{.30\linewidth}c p{.55\linewidth}@{}}\toprule
Complete outcomes or categories & $\GTenHLValue{x}$ (\GTenHLValue{USD}) & Grouped probability: $\GTenHLProbability{p_X(}\GTenHLValue{x}\GTenHLProbability{)}$\\\midrule
$\GTenHLOutcome{(3,0,0)}$: \GTenHLSecond{1} subsets & $\GTenHLValue{0}$ & $\begin{aligned}[t]\GTenHLProbability{p_X(}\GTenHLValue{0}\GTenHLProbability{)}&=\GTenHLProbability{1/84}\\&\approx\GTenHLProbability{0.0119}=\GTenHLProbability{1.19\%}\end{aligned}$\\[5pt]
$\GTenHLOutcome{(2,1,0)}$: \GTenHLSecond{12} subsets & $\GTenHLValue{6}$ & $\begin{aligned}[t]\GTenHLProbability{p_X(}\GTenHLValue{6}\GTenHLProbability{)}&=\GTenHLProbability{12/84}=\GTenHLProbability{1/7}\\&\approx\GTenHLProbability{0.1429}=\GTenHLProbability{14.29\%}\end{aligned}$\\[5pt]
$\GTenHLOutcome{(1,2,0)}$: \GTenHLSecond{18} subsets & $\GTenHLValue{12}$ & $\begin{aligned}[t]\GTenHLProbability{p_X(}\GTenHLValue{12}\GTenHLProbability{)}&=\GTenHLProbability{18/84}=\GTenHLProbability{3/14}\\&\approx\GTenHLProbability{0.2143}=\GTenHLProbability{21.43\%}\end{aligned}$\\[5pt]
$\GTenHLOutcome{(2,0,1)}$: \GTenHLSecond{6} subsets & $\GTenHLValue{14}$ & $\begin{aligned}[t]\GTenHLProbability{p_X(}\GTenHLValue{14}\GTenHLProbability{)}&=\GTenHLProbability{6/84}=\GTenHLProbability{1/14}\\&\approx\GTenHLProbability{0.0714}=\GTenHLProbability{7.14\%}\end{aligned}$\\[5pt]
$\GTenHLOutcome{(0,3,0)}$: \GTenHLSecond{4} subsets & $\GTenHLValue{18}$ & $\begin{aligned}[t]\GTenHLProbability{p_X(}\GTenHLValue{18}\GTenHLProbability{)}&=\GTenHLProbability{4/84}=\GTenHLProbability{1/21}\\&\approx\GTenHLProbability{0.0476}=\GTenHLProbability{4.76\%}\end{aligned}$\\[5pt]
$\GTenHLOutcome{(1,1,1)}$: \GTenHLSecond{24} subsets\newline $\GTenHLOutcome{(1,0,2)}$: \GTenHLSecond{3} subsets & $\GTenHLValue{28}$ & $\begin{aligned}[t]\GTenHLProbability{p_X(}\GTenHLValue{28}\GTenHLProbability{)}&=\GTenHLProbability{24/84}+\GTenHLProbability{3/84}=\GTenHLProbability{27/84}=\GTenHLProbability{9/28}\\&\approx\GTenHLProbability{0.3214}=\GTenHLProbability{32.14\%}\end{aligned}$\\[5pt]
$\GTenHLOutcome{(0,2,1)}$: \GTenHLSecond{12} subsets\newline $\GTenHLOutcome{(0,1,2)}$: \GTenHLSecond{4} subsets & $\GTenHLValue{30}$ & $\begin{aligned}[t]\GTenHLProbability{p_X(}\GTenHLValue{30}\GTenHLProbability{)}&=\GTenHLProbability{12/84}+\GTenHLProbability{4/84}=\GTenHLProbability{16/84}=\GTenHLProbability{4/21}\\&\approx\GTenHLProbability{0.1905}=\GTenHLProbability{19.05\%}\end{aligned}$\\[5pt]
\bottomrule\end{tabular}
\endgroup
\end{center}
\begin{center}\begin{tikzpicture}
\begin{axis}[keyaxis,width=.86\linewidth,height=5cm,xmin=-2,xmax=32,ymin=0,ymax=.65,xtick={0,6,12,14,18,28,30},xticklabel style={rotate=45,anchor=north east},ylabel={$\GTenHLProbability{p_X(}\GTenHLValue{x}\GTenHLProbability{)}$},yticklabel={$\GTenHLProbability{\pgfmathprintnumber{\tick}}$}]
\draw[fill=keyMassFill,draw=keyMassColor] (axis cs:-0.225,0) rectangle (axis cs:0.225,1/84);
\draw[fill=keyMassFill,draw=keyMassColor] (axis cs:5.775,0) rectangle (axis cs:6.225,12/84);
\draw[fill=keyMassFill,draw=keyMassColor] (axis cs:11.775,0) rectangle (axis cs:12.225,18/84);
\draw[fill=keyMassFill,draw=keyMassColor] (axis cs:13.775,0) rectangle (axis cs:14.225,6/84);
\draw[fill=keyMassFill,draw=keyMassColor] (axis cs:17.775,0) rectangle (axis cs:18.225,4/84);
\draw[fill=keyMassFill,draw=keyMassColor] (axis cs:27.775,0) rectangle (axis cs:28.225,27/84);
\draw[fill=keyMassFill,draw=keyMassColor] (axis cs:29.775,0) rectangle (axis cs:30.225,16/84);
\end{axis}\end{tikzpicture}\end{center}\Needspace{5\baselineskip}\textbf{PMF checks.}\begin{itemize}\item Every mass is positive.\item Normalization: $(1+12+18+6+4+27+16)/84=1$.\item Set $\GTenHLProbability{p_X(}\GTenHLValue{x}\GTenHLProbability{)}=\GTenHLProbability{0}$ outside the support. Bars use actual monetary spacing; gaps carry no mass.\end{itemize}

\textbf{Scoring note.} These explanatory positivity and normalization checks carry no separate points; omitting them does not reduce the task score.

\par\smallskip
\noindent\begin{minipage}{\linewidth}\centering
\textbf{Proposed partial credit.}\par\smallskip
\begin{tabular}{@{}l r@{}}\toprule
\CreditDescription{\textbf{Assessed step or criterion}} & \textbf{Points}\\\midrule
\CreditDescription{Variable, units and support} & 0.125\\
\CreditDescription{The denominator and exhaustive composition counts} & 0.5\\
\CreditDescription{The bonus, cap and aggregation} & 0.375\\
\CreditDescription{The grouped PMF table} & 0.25\\
\CreditDescription{The labeled, correctly spaced graph} & 0.25\\
\midrule
\CreditDescription{\textbf{Total}} & \textbf{1.5}\\
\bottomrule\end{tabular}
\end{minipage}\par\smallskip

\par\vspace{3mm}\secInicio[][sec:q6]

\clearpage
\textbf{Question 6 --- continued}\quad\hyperref[sec:q6]{Original description}

\Needspace{8\baselineskip}
\subsection*{Task 2 --- question and worked solution (proposed allocation) \score{3}}
\label{sol:q6t2}

\noindent\textbf{\boldmath Construct the complete CDF table, including both tails, and draw a labeled step graph. Use open and filled endpoints correctly and preserve actual monetary spacing.}\par\nopagebreak

For $\HLCumulativeProbability{F_X(}\GTenHLValue{x}\HLCumulativeProbability{)}=P(X\leq x)$, add the masses in increasing monetary order.

The successive cumulative numerators are $1,13,31,37,41,68,84$, all over $84$.
\begin{center}
\begingroup
\setlength{\tabcolsep}{5pt}
\begin{tabular}{@{}r@{}c@{}c@{}c@{}llc@{}}\toprule
\multicolumn{5}{c}{Cutoff interval (USD)} & Calculations & $\HLCumulativeProbability{F_X(}\GTenHLValue{x}\HLCumulativeProbability{)}$\\\midrule
 & & $\GTenHLValue{x}$ & $<$ & $\GTenHLValue{0}$ & $\text{No included support values}$ & $\HLCumulativeProbability{0}$\\[5pt]
$\GTenHLValue{0}$ & $\leq$ & $\GTenHLValue{x}$ & $<$ & $\GTenHLValue{6}$ & $\begin{aligned}[t]
  &\GTenHLProbability{p_X(}\GTenHLValue{0}\GTenHLProbability{)}
  \\&= \GTenHLProbability{\frac{1}{84}}
\end{aligned}$ & $\HLCumulativeProbability{1/84}$\\[5pt]
$\GTenHLValue{6}$ & $\leq$ & $\GTenHLValue{x}$ & $<$ & $\GTenHLValue{12}$ & $\begin{aligned}[t]
  &\GTenHLProbability{p_X(}\GTenHLValue{0}\GTenHLProbability{)} + \GTenHLProbability{p_X(}\GTenHLValue{6}\GTenHLProbability{)}
  \\&= \GTenHLProbability{\frac{1}{84}} + \GTenHLProbability{\frac{12}{84}}
\end{aligned}$ & $\HLCumulativeProbability{13/84}$\\[5pt]
$\GTenHLValue{12}$ & $\leq$ & $\GTenHLValue{x}$ & $<$ & $\GTenHLValue{14}$ & $\begin{aligned}[t]
  &\GTenHLProbability{p_X(}\GTenHLValue{0}\GTenHLProbability{)} + \GTenHLProbability{p_X(}\GTenHLValue{6}\GTenHLProbability{)} + \GTenHLProbability{p_X(}\GTenHLValue{12}\GTenHLProbability{)}
  \\&= \GTenHLProbability{\frac{1}{84}} + \GTenHLProbability{\frac{12}{84}} + \GTenHLProbability{\frac{18}{84}}
\end{aligned}$ & $\HLCumulativeProbability{31/84}$\\[5pt]
$\GTenHLValue{14}$ & $\leq$ & $\GTenHLValue{x}$ & $<$ & $\GTenHLValue{18}$ & $\begin{aligned}[t]
  &\GTenHLProbability{p_X(}\GTenHLValue{0}\GTenHLProbability{)} + \GTenHLProbability{p_X(}\GTenHLValue{6}\GTenHLProbability{)} + \GTenHLProbability{p_X(}\GTenHLValue{12}\GTenHLProbability{)} + \GTenHLProbability{p_X(}\GTenHLValue{14}\GTenHLProbability{)}
  \\&= \GTenHLProbability{\frac{1}{84}} + \GTenHLProbability{\frac{12}{84}} + \GTenHLProbability{\frac{18}{84}} + \GTenHLProbability{\frac{6}{84}}
\end{aligned}$ & $\HLCumulativeProbability{37/84}$\\[5pt]
$\GTenHLValue{18}$ & $\leq$ & $\GTenHLValue{x}$ & $<$ & $\GTenHLValue{28}$ & $\begin{aligned}[t]
  &\GTenHLProbability{p_X(}\GTenHLValue{0}\GTenHLProbability{)} + \GTenHLProbability{p_X(}\GTenHLValue{6}\GTenHLProbability{)} + \GTenHLProbability{p_X(}\GTenHLValue{12}\GTenHLProbability{)} + \GTenHLProbability{p_X(}\GTenHLValue{14}\GTenHLProbability{)}
  \\&{}+ \GTenHLProbability{p_X(}\GTenHLValue{18}\GTenHLProbability{)}
  \\&= \GTenHLProbability{\frac{1}{84}} + \GTenHLProbability{\frac{12}{84}} + \GTenHLProbability{\frac{18}{84}} + \GTenHLProbability{\frac{6}{84}}
  \\&{}+ \GTenHLProbability{\frac{4}{84}}
\end{aligned}$ & $\HLCumulativeProbability{41/84}$\\[5pt]
$\GTenHLValue{28}$ & $\leq$ & $\GTenHLValue{x}$ & $<$ & $\GTenHLValue{30}$ & $\begin{aligned}[t]
  &\GTenHLProbability{p_X(}\GTenHLValue{0}\GTenHLProbability{)} + \GTenHLProbability{p_X(}\GTenHLValue{6}\GTenHLProbability{)} + \GTenHLProbability{p_X(}\GTenHLValue{12}\GTenHLProbability{)} + \GTenHLProbability{p_X(}\GTenHLValue{14}\GTenHLProbability{)}
  \\&{}+ \GTenHLProbability{p_X(}\GTenHLValue{18}\GTenHLProbability{)} + \GTenHLProbability{p_X(}\GTenHLValue{28}\GTenHLProbability{)}
  \\&= \GTenHLProbability{\frac{1}{84}} + \GTenHLProbability{\frac{12}{84}} + \GTenHLProbability{\frac{18}{84}} + \GTenHLProbability{\frac{6}{84}}
  \\&{}+ \GTenHLProbability{\frac{4}{84}} + \GTenHLProbability{\frac{27}{84}}
\end{aligned}$ & $\HLCumulativeProbability{17/21}$\\[5pt]
$\GTenHLValue{30}$ & $\leq$ & $\GTenHLValue{x}$ &  &  & $\begin{aligned}[t]
  &\GTenHLProbability{p_X(}\GTenHLValue{0}\GTenHLProbability{)} + \GTenHLProbability{p_X(}\GTenHLValue{6}\GTenHLProbability{)} + \GTenHLProbability{p_X(}\GTenHLValue{12}\GTenHLProbability{)} + \GTenHLProbability{p_X(}\GTenHLValue{14}\GTenHLProbability{)}
  \\&{}+ \GTenHLProbability{p_X(}\GTenHLValue{18}\GTenHLProbability{)} + \GTenHLProbability{p_X(}\GTenHLValue{28}\GTenHLProbability{)} + \GTenHLProbability{p_X(}\GTenHLValue{30}\GTenHLProbability{)}
  \\&= \GTenHLProbability{\frac{1}{84}} + \GTenHLProbability{\frac{12}{84}} + \GTenHLProbability{\frac{18}{84}} + \GTenHLProbability{\frac{6}{84}}
  \\&{}+ \GTenHLProbability{\frac{4}{84}} + \GTenHLProbability{\frac{27}{84}} + \GTenHLProbability{\frac{16}{84}}
\end{aligned}$ & $\HLCumulativeProbability{1}$\\[5pt]
\bottomrule\end{tabular}
\endgroup
\end{center}
\begin{center}\begin{tikzpicture}
\begin{axis}[keyaxis,width=.86\linewidth,height=7cm,xmin=-2,xmax=32,ymin=-.07,ymax=1.1,xtick={0,6,12,14,18,28,30},xticklabel style={rotate=45,anchor=north east},ylabel={$\HLCumulativeProbability{F_X(}\GTenHLValue{x}\HLCumulativeProbability{)}$},yticklabel={$\HLCumulativeProbability{\pgfmathprintnumber{\tick}}$},ytick={0,.25,.5,.75,1}]
\addplot[keyCDFColor,thick,<-] coordinates {(-2,0)(0,0)};
\addplot[keyCDFColor,thick] coordinates {(0,1/84)(6,1/84)};
\addplot[keyCDFColor,thick] coordinates {(6,13/84)(12,13/84)};
\addplot[keyCDFColor,thick] coordinates {(12,31/84)(14,31/84)};
\addplot[keyCDFColor,thick] coordinates {(14,37/84)(18,37/84)};
\addplot[keyCDFColor,thick] coordinates {(18,41/84)(28,41/84)};
\addplot[keyCDFColor,thick] coordinates {(28,68/84)(30,68/84)};
\addplot[keyCDFColor,thick,->] coordinates {(30,1)(32,1)};
\addplot[keyCDFColor,only marks,mark=*,mark size=.8pt,mark options={fill=white}] coordinates {(0,0/84)(6,1/84)(12,13/84)(14,31/84)(18,37/84)(28,41/84)(30,68/84)};
\addplot[keyCDFColor,only marks,mark=*,mark size=.8pt] coordinates {(0,1/84)(6,13/84)(12,31/84)(14,37/84)(18,41/84)(28,68/84)(30,84/84)};
\end{axis}\end{tikzpicture}\end{center}
\clearpage
\textbf{Question 6 --- Task 2 continued: comparison}\quad\hyperref[sec:q6]{Original description}
\par\smallskip\textbf{Same distribution: CDF above, PMF below.}
\begin{center}\begin{minipage}{\linewidth}\centering
\begin{tikzpicture}
\pgfplotsset{keySupport/.style={xmin=-2,xmax=32,xtick={0,6,12,14,18,28,30},
xlabel={$\GTenHLValue{x}$ (\GTenHLValue{USD})},
xticklabel={$\GTenHLValue{\pgfmathprintnumber{\tick}}$},
xticklabel style={font=\footnotesize,anchor=north,fill=white,inner sep=1pt},
xlabel style={fill=white,inner sep=2pt}}}
\begin{axis}[keyCombined,keySupport,name=keyTop,anchor=south west,
ylabel={$\HLCumulativeProbability{F_X(}\GTenHLValue{x}\HLCumulativeProbability{)}$},yticklabel={$\HLCumulativeProbability{\pgfmathprintnumber{\tick}}$},keyIncluded/.append style={mark size=.5pt},keyExcluded/.append style={mark size=.5pt}]
\keyPhantom{0}{1/84}{0/84}
\keyPhantom{6}{12/84}{1/84}
\keyPhantom{12}{18/84}{13/84}
\keyPhantom{14}{6/84}{31/84}
\keyPhantom{18}{4/84}{37/84}
\keyPhantom{28}{27/84}{41/84}
\keyPhantom{30}{16/84}{68/84}
\addplot[keyStep,<-] coordinates {(-2,0)(0,0)};
\addplot[keyStep] coordinates {(0,1/84)(6,1/84)};
\addplot[keyStep] coordinates {(6,13/84)(12,13/84)};
\addplot[keyStep] coordinates {(12,31/84)(14,31/84)};
\addplot[keyStep] coordinates {(14,37/84)(18,37/84)};
\addplot[keyStep] coordinates {(18,41/84)(28,41/84)};
\addplot[keyStep] coordinates {(28,68/84)(30,68/84)};
\addplot[keyStep,->] coordinates {(30,1)(32,1)};
\addplot[keyExcluded] coordinates {(0,0/84)(6,1/84)(12,13/84)(14,31/84)(18,37/84)(28,41/84)(30,68/84)};
\addplot[keyIncluded] coordinates {(0,1/84)(6,13/84)(12,31/84)(14,37/84)(18,41/84)(28,68/84)(30,84/84)};
\end{axis}
\begin{axis}[keyCombined,keySupport,at={(keyTop.south west)},
anchor=north west,yshift=-1.35cm,ylabel={$\GTenHLProbability{p_X(}\GTenHLValue{x}\GTenHLProbability{)}$},
yticklabel={$\GTenHLProbability{\pgfmathprintnumber{\tick}}$}]
\keyMass{0}{1/84}
\keyMass{6}{12/84}
\keyMass{12}{18/84}
\keyMass{14}{6/84}
\keyMass{18}{4/84}
\keyMass{28}{27/84}
\keyMass{30}{16/84}
\addlegendimage{keyStep}\addlegendentry{Top: \HLCumulativeProbability{CDF} steps}
\addlegendimage{keyPhantom,mark=*}\addlegendentry{Top: \GTenHLProbability{PMF} shadows (jumps)}
\addlegendimage{keyMass,mark=*}\addlegendentry{Bottom: \GTenHLProbability{PMF} stems}
\addlegendimage{keyCorrespondence}\addlegendentry{Midpoint correspondence}
\end{axis}
\keyCorrespondence{0}
\keyCorrespondence{6}
\keyCorrespondence{12}
\keyCorrespondence{14}
\keyCorrespondence{18}
\keyCorrespondence{28}
\keyCorrespondence{30}
\end{tikzpicture}\end{minipage}\end{center}
\keyConnectionReading
\par\medskip
\textbf{CDF checks.}\begin{itemize}\item The intervals cover every real cutoff, and every height lies in $[0,1]$.\item The heights never decrease.\item Each segment includes its left endpoint and excludes its right endpoint; the filled point gives the value at a jump, so the CDF is right-continuous.\item The lower-tail limit is 0 and the upper-tail limit is 1; the arrows extend both tails.\item Successive cumulative differences recover the PMF numerators $1,12,18,6,4,27,16$ over $84$. Thus $\HLCumulativeProbability{F_X(}\GTenHLValue{a}\HLCumulativeProbability{)}-\HLCumulativeProbability{F_X(}\GTenHLValue{a^-}\HLCumulativeProbability{)}=\GTenHLProbability{p_X(}\GTenHLValue{a}\GTenHLProbability{)}$; unsupported cutoffs have zero jump and zero mass.\end{itemize}

\par\smallskip
\noindent\begin{minipage}{\linewidth}\centering
\textbf{Proposed partial credit.}\par\smallskip
\begin{tabular}{@{}l r@{}}\toprule
\CreditDescription{\textbf{Assessed step or criterion}} & \textbf{Points}\\\midrule
\CreditDescription{Cumulative sums} & 0.75\\
\CreditDescription{The complete interval table and tails} & 0.75\\
\CreditDescription{The labeled, correctly spaced step graph and endpoints} & 0.75\\
\CreditDescription{The CDF and jump checks} & 0.75\\
\midrule
\CreditDescription{\textbf{Total}} & \textbf{3}\\
\bottomrule\end{tabular}
\end{minipage}\par\smallskip

\par\vspace{3mm}\secInicio[][sec:q6]

\clearpage
\textbf{Question 6 --- continued}\quad\hyperref[sec:q6]{Original description}

\Needspace{8\baselineskip}
\subsection*{Task 3 (c) --- question and worked solution (proposed allocation) \score{1.5}}
\label{sol:q6t3}

\noindent\textbf{\boldmath Calculate and interpret (i) the probability that the final payout is less than 14 USD and (ii) the probability that it is at least 28 USD.}\par\nopagebreak

\textbf{Allocation:} (i) 0.75 points; (ii) 0.75 points.

\textbf{(i)} Include only $0,6,12$ USD; exclude the boundary mass at 14 USD.
\[
\begin{aligned}
P(\GTenHLValue{X}<\GTenHLValue{14})&=\GTenHLProbability{p_X(}\GTenHLValue{0}\GTenHLProbability{)}+\GTenHLProbability{p_X(}\GTenHLValue{6}\GTenHLProbability{)}+\GTenHLProbability{p_X(}\GTenHLValue{12}\GTenHLProbability{)}=\frac{1+12+18}{84}
\\ &=\GTenHLProbability{\frac{31}{84}}=\GTenHLProbability{0.36\overline{904761}}=\GTenHLProbability{36.\overline{904761}\%}.
\end{aligned}
\]
The final payout for a uniformly selected three-ticket subset, after the bonus and cap is below 14 USD with probability $\GTenHLProbability{31/84}\approx\GTenHLProbability{36.90\%}$.

\textbf{(ii)} Include only $28,30$ USD; include the boundary mass at 28 USD.
\[
\begin{aligned}
P(\GTenHLValue{X}\geq\GTenHLValue{28})&=\GTenHLProbability{p_X(}\GTenHLValue{28}\GTenHLProbability{)}+\GTenHLProbability{p_X(}\GTenHLValue{30}\GTenHLProbability{)}=\frac{27+16}{84}
\\ &=\GTenHLProbability{\frac{43}{84}}=\GTenHLProbability{0.51\overline{190476}}=\GTenHLProbability{51.\overline{190476}\%}.
\end{aligned}
\]
The final payout for a uniformly selected three-ticket subset, after the bonus and cap is at least 28 USD with probability $\GTenHLProbability{43/84}\approx\GTenHLProbability{51.19\%}$.

\par\smallskip
\noindent\begin{minipage}{\linewidth}\centering
\textbf{Proposed partial credit.}\par\smallskip
\begin{tabular}{@{}l r@{}}\toprule
\CreditDescription{\textbf{Assessed step or criterion}} & \textbf{Points}\\\midrule
\CreditDescription{(i): selecting exactly the included PMF masses with the correct boundary} & 0.375\\
\CreditDescription{(i): summing them correctly} & 0.1875\\
\CreditDescription{(i): the contextual interpretation} & 0.1875\\
\CreditDescription{(ii): selecting exactly the included PMF masses with the correct boundary} & 0.375\\
\CreditDescription{(ii): summing them correctly} & 0.1875\\
\CreditDescription{(ii): the contextual interpretation} & 0.1875\\
\midrule
\CreditDescription{\textbf{Total}} & \textbf{1.5}\\
\bottomrule\end{tabular}
\end{minipage}\par\smallskip

Each part totals 0.75 points.

\par\vspace{3mm}\secInicio[][sec:q6]

\textbf{Scoring reconciliation.} Question 6 gives $1.5+3+1.5=6$ points for C2. C2 gives $2+2+4+6=14$ points; its independent passing threshold is 4 points. C1 gives $3+3=6$ points; its independent passing threshold is also 4 points. The overall available total is $6+14=20$ points, with no transfer between criteria.

\end{document}
